
\documentclass{article}
\usepackage{colortbl}
\usepackage{makecell}
\usepackage{multirow}
\usepackage{supertabular}

\begin{document}

\newcounter{utterance}

\twocolumn

{ \footnotesize  \setcounter{utterance}{1}
\setlength{\tabcolsep}{0pt}
\begin{supertabular}{c@{$\;$}|p{.15\linewidth}@{}p{.15\linewidth}p{.15\linewidth}p{.15\linewidth}p{.15\linewidth}p{.15\linewidth}}

    \# & $\;$A & \multicolumn{4}{c}{Game Master} & $\;\:$B\\
    \hline 

    \theutterance \stepcounter{utterance}  

    & & \multicolumn{4}{p{0.6\linewidth}}{\cellcolor[rgb]{0.9,0.9,0.9}{%
	\makecell[{{p{\linewidth}}}]{% 
	  \tt {\tiny [A$\langle$GM]}  
	 Sono il tuo game master, e stai giocando un gioco collaborativo con la seguente griglia come tavolo da gioco:\\ \tt \\ \tt ```\\ \tt  1234567\\ \tt ╔═╤═╤═╤═╗\\ \tt ║◌│◌│◌│◌║ 1\\ \tt ║P└─┼─┴─╢ 2\\ \tt ║◌SA│◌◌◌║ 3\\ \tt ╟──┬┴───╢ 4\\ \tt ║◌◌│◌◌O◌║ 5\\ \tt ║◌◌│◌◌MD║ 6\\ \tt ║◌◌│◌◌◌T║ 7\\ \tt ╚══╧════╝\\ \tt \\ \tt ```\\ \tt \\ \tt * Il bordo superiore mostra le coordinate x che aumentano verso destra, e il bordo destro le coordinate y che aumentano verso il basso.\\ \tt * The following objects are randomly placed on your grid: 'P', 'O', 'T', 'S', 'D', 'A', 'M'.\\ \tt \\ \tt \\ \tt L'altro giocatore vede una variante del tavolo da gioco, dove gli oggetti sono posizionati in diverse posizioni casuali. Non puoi vedere il tavolo dell'altro giocatore, e loro non possono vedere il tuo.\\ \tt \\ \tt \\ \tt **Obiettivo:**\\ \tt \\ \tt Entrambi i giocatori devono spostare gli oggetti sui rispettivi sfondi in modo che oggetti identici finiscano alle stesse coordinate. Devi comunicare con l'altro giocatore per concordare una configurazione obiettivo comune.\\ \tt \\ \tt \\ \tt **Regole:**\\ \tt \\ \tt * In ogni turno, puoi inviare esattamente uno dei seguenti due comandi:\\ \tt  1. `COMUNICA: <MESSAGE>`: per inviare un messaggio (tutto fino alla prossima interruzione di riga) all'altro giocatore. Lo inoltrerò al tuo partner.\\ \tt 2. `SPOSTA: <OBJECT>, (<X>, <Y>)`: per spostare un oggetto in una nuova posizione, dove <X> è la colonna e <Y> è la riga. Ti farò sapere se la tua mossa è andata a buon fine o no.\\ \tt * Se non rispetti il formato, o invii diversi comandi contemporaneamente, devo penalizzarti.\\ \tt * If both players accumulate more than 16 penalties, you both lose the game.\\ \tt * È vitale che comunichi con l'altro giocatore riguardo al tuo stato obiettivo! L'*unico* modo in cui puoi trasmettere la tua strategia all'altro giocatore è usando il comando `COMUNICA: <MESSAGE>`!\\ \tt \\ \tt \\ \tt **Spostamento Oggetti**\\ \tt \\ \tt * Puoi spostare oggetti solo in celle entro i limiti della griglia. La cella di destinazione deve essere vuota, cioè deve contenere solo il simbolo '◌'.\\ \tt * Se provi a spostare un oggetto in un posto che non è vuoto, o provi a spostarlo fuori dalla griglia, devo penalizzarti. Hai un altro tentativo.\\ \tt * Prima di fare una mossa, controlla due volte che il posto di destinazione sia vuoto, e non contenga alcuna lettera, cornice, o linea!\\ \tt \\ \tt \\ \tt **Fine del Gioco**\\ \tt \\ \tt Se pensi di aver raggiunto l'obiettivo di allineare tutti gli oggetti, puoi chiedere all'altro giocatore di finire il gioco inviando COMUNICA: finito?. Se l'altro giocatore ti chiede di finire il gioco, e rispondi con COMUNICA: finito!, the game will end.\\ \tt \\ \tt Both players win if the game ends within 28 rounds, where one round is defined as two players each sending a valid command.\\ \tt \\ \tt **Punteggio:**\\ \tt \\ \tt Più vicini sono gli oggetti identici in entrambi i tavoli da gioco, più punti ottieni. Le penalità riducono i tuoi punti. Riesci a battere il record?\\ \tt \\ \tt \\ \tt Per favore invia un messaggio all'altro giocatore per iniziare il gioco! 
	  } 
	   } 
	   } 
	 & \\ 
 

    \theutterance \stepcounter{utterance}  

    & \multicolumn{4}{p{0.6\linewidth}}{\cellcolor[rgb]{0.8,1,0.9}{%
	\makecell[{{p{\linewidth}}}]{% 
	  \tt {\tiny [A$\rangle$GM]}  
	 COMUNICA: Ciao! Sono pronto a iniziare. Propongo di iniziare con un oggetto facile come "P". Dove si trova il tuo "P"? 
	  } 
	   } 
	   } 
	 & & \\ 
 

    \theutterance \stepcounter{utterance}  

    & & & \multicolumn{2}{p{0.3\linewidth}}{\cellcolor[rgb]{0.95,0.95,0.95}{%
	\makecell[{{p{\linewidth}}}]{% 
	  \tt {\tiny [GM$|$GM]}  
	 COMUNICA: Ciao! Sono pronto a iniziare. Propongo di iniziare con un oggetto facile come "P". Dove si trova il tuo "P"? 
	  } 
	   } 
	   } 
	 & & \\ 
 

    \theutterance \stepcounter{utterance}  

    & & \multicolumn{4}{p{0.6\linewidth}}{\cellcolor[rgb]{0.9,0.9,0.9}{%
	\makecell[{{p{\linewidth}}}]{% 
	  \tt {\tiny [GM$\rangle$B]}  
	 Sono il tuo game master, e stai giocando un gioco collaborativo con la seguente griglia come tavolo da gioco:\\ \tt \\ \tt ```\\ \tt  1234567\\ \tt ╔═╤═╤═╤═╗\\ \tt ║◌│◌│◌│P║ 1\\ \tt ║T└─┼─┴─╢ 2\\ \tt ║◌◌◌│M◌◌║ 3\\ \tt ╟──┬┴───╢ 4\\ \tt ║◌D│◌◌O◌║ 5\\ \tt ║◌◌│◌S◌◌║ 6\\ \tt ║A◌│◌◌◌◌║ 7\\ \tt ╚══╧════╝\\ \tt \\ \tt ```\\ \tt \\ \tt * Il bordo superiore mostra le coordinate x che aumentano verso destra, e il bordo destro le coordinate y che aumentano verso il basso.\\ \tt * The following objects are randomly placed on your grid: 'P', 'O', 'T', 'S', 'D', 'A', 'M'.\\ \tt \\ \tt \\ \tt L'altro giocatore vede una variante del tavolo da gioco, dove gli oggetti sono posizionati in diverse posizioni casuali. Non puoi vedere il tavolo dell'altro giocatore, e loro non possono vedere il tuo.\\ \tt \\ \tt \\ \tt **Obiettivo:**\\ \tt \\ \tt Entrambi i giocatori devono spostare gli oggetti sui rispettivi sfondi in modo che oggetti identici finiscano alle stesse coordinate. Devi comunicare con l'altro giocatore per concordare una configurazione obiettivo comune.\\ \tt \\ \tt \\ \tt **Regole:**\\ \tt \\ \tt * In ogni turno, puoi inviare esattamente uno dei seguenti due comandi:\\ \tt  1. `COMUNICA: <MESSAGE>`: per inviare un messaggio (tutto fino alla prossima interruzione di riga) all'altro giocatore. Lo inoltrerò al tuo partner.\\ \tt 2. `SPOSTA: <OBJECT>, (<X>, <Y>)`: per spostare un oggetto in una nuova posizione, dove <X> è la colonna e <Y> è la riga. Ti farò sapere se la tua mossa è andata a buon fine o no.\\ \tt * Se non rispetti il formato, o invii diversi comandi contemporaneamente, devo penalizzarti.\\ \tt * If both players accumulate more than 16 penalties, you both lose the game.\\ \tt * È vitale che comunichi con l'altro giocatore riguardo al tuo stato obiettivo! L'*unico* modo in cui puoi trasmettere la tua strategia all'altro giocatore è usando il comando `COMUNICA: <MESSAGE>`!\\ \tt \\ \tt \\ \tt **Spostamento Oggetti**\\ \tt \\ \tt * Puoi spostare oggetti solo in celle entro i limiti della griglia. La cella di destinazione deve essere vuota, cioè deve contenere solo il simbolo '◌'.\\ \tt * Se provi a spostare un oggetto in un posto che non è vuoto, o provi a spostarlo fuori dalla griglia, devo penalizzarti. Hai un altro tentativo.\\ \tt * Prima di fare una mossa, controlla due volte che il posto di destinazione sia vuoto, e non contenga alcuna lettera, cornice, o linea!\\ \tt \\ \tt \\ \tt **Fine del Gioco**\\ \tt \\ \tt Se pensi di aver raggiunto l'obiettivo di allineare tutti gli oggetti, puoi chiedere all'altro giocatore di finire il gioco inviando COMUNICA: finito?. Se l'altro giocatore ti chiede di finire il gioco, e rispondi con COMUNICA: finito!, the game will end.\\ \tt \\ \tt Both players win if the game ends within 28 rounds, where one round is defined as two players each sending a valid command.\\ \tt \\ \tt **Punteggio:**\\ \tt \\ \tt Più vicini sono gli oggetti identici in entrambi i tavoli da gioco, più punti ottieni. Le penalità riducono i tuoi punti. Riesci a battere il record?\\ \tt \\ \tt L'altro giocatore ha iniziato il gioco inviando questo messaggio:\\ \tt \\ \tt "Ciao! Sono pronto a iniziare. Propongo di iniziare con un oggetto facile come "P". Dove si trova il tuo "P"?"\\ \tt \\ \tt Qual è il tuo primo comando? 
	  } 
	   } 
	   } 
	 & \\ 
 

    \theutterance \stepcounter{utterance}  

    & & & \multicolumn{4}{p{0.6\linewidth}}{\cellcolor[rgb]{1,0.85,0.72}{%
	\makecell[{{p{\linewidth}}}]{% 
	  \tt {\tiny [GM$\langle$B]}  
	 COMUNICA: Ciao! Il mio "P" si trova alla posizione (4,1). E il tuo? Dove si trova il tuo "P"? 
	  } 
	   } 
	   } 
	  \\ 
 

    \theutterance \stepcounter{utterance}  

    & & & \multicolumn{2}{p{0.3\linewidth}}{\cellcolor[rgb]{0.95,0.95,0.95}{%
	\makecell[{{p{\linewidth}}}]{% 
	  \tt {\tiny [GM$|$GM]}  
	 COMUNICA: Ciao! Il mio "P" si trova alla posizione (4,1). E il tuo? Dove si trova il tuo "P"? 
	  } 
	   } 
	   } 
	 & & \\ 
 

    \theutterance \stepcounter{utterance}  

    & & \multicolumn{4}{p{0.6\linewidth}}{\cellcolor[rgb]{0.9,0.9,0.9}{%
	\makecell[{{p{\linewidth}}}]{% 
	  \tt {\tiny [A$\langle$GM]}  
	 Il tuo messaggio è stato inoltrato all'altro giocatore.\\ \tt \\ \tt \\ \tt You are currently playing round 1 of maximum 28.\\ \tt You have collectively accumulated 0 of 16 penalties.\\ \tt L'altro giocatore ha inviato questo messaggio:\\ \tt \\ \tt \\ \tt "Ciao! Il mio "P" si trova alla posizione (4,1). E il tuo? Dove si trova il tuo "P"?"\\ \tt \\ \tt \\ \tt Qual è il tuo prossimo comando? 
	  } 
	   } 
	   } 
	 & \\ 
 

    \theutterance \stepcounter{utterance}  

    & \multicolumn{4}{p{0.6\linewidth}}{\cellcolor[rgb]{0.8,1,0.9}{%
	\makecell[{{p{\linewidth}}}]{% 
	  \tt {\tiny [A$\rangle$GM]}  
	 COMUNICA: Il mio "P" si trova alla posizione (1,2). Interessante, il tuo "P" è più a destra del mio. Propongo di spostare i nostri "P" in una posizione centrale, ad esempio (3,2) o (3,3), per vedere se possiamo allineare gli altri oggetti intorno a esso. Che ne pensi? 
	  } 
	   } 
	   } 
	 & & \\ 
 

    \theutterance \stepcounter{utterance}  

    & & & \multicolumn{2}{p{0.3\linewidth}}{\cellcolor[rgb]{0.95,0.95,0.95}{%
	\makecell[{{p{\linewidth}}}]{% 
	  \tt {\tiny [GM$|$GM]}  
	 COMUNICA: Il mio "P" si trova alla posizione (1,2). Interessante, il tuo "P" è più a destra del mio. Propongo di spostare i nostri "P" in una posizione centrale, ad esempio (3,2) o (3,3), per vedere se possiamo allineare gli altri oggetti intorno a esso. Che ne pensi? 
	  } 
	   } 
	   } 
	 & & \\ 
 

    \theutterance \stepcounter{utterance}  

    & & \multicolumn{4}{p{0.6\linewidth}}{\cellcolor[rgb]{0.9,0.9,0.9}{%
	\makecell[{{p{\linewidth}}}]{% 
	  \tt {\tiny [GM$\rangle$B]}  
	 Il tuo messaggio è stato inoltrato all'altro giocatore.\\ \tt \\ \tt \\ \tt You are currently playing round 2 of maximum 28.\\ \tt You have collectively accumulated 0 of 16 penalties.\\ \tt L'altro giocatore ha inviato questo messaggio:\\ \tt \\ \tt \\ \tt "Il mio "P" si trova alla posizione (1,2). Interessante, il tuo "P" è più a destra del mio. Propongo di spostare i nostri "P" in una posizione centrale, ad esempio (3,2) o (3,3), per vedere se possiamo allineare gli altri oggetti intorno a esso. Che ne pensi?"\\ \tt \\ \tt \\ \tt Qual è il tuo prossimo comando? 
	  } 
	   } 
	   } 
	 & \\ 
 

    \theutterance \stepcounter{utterance}  

    & & & \multicolumn{4}{p{0.6\linewidth}}{\cellcolor[rgb]{1,0.85,0.72}{%
	\makecell[{{p{\linewidth}}}]{% 
	  \tt {\tiny [GM$\langle$B]}  
	 COMUNICA: Sono d'accordo, spostare i "P" in una posizione centrale potrebbe essere un buon inizio. Tuttavia, prima di prendere una decisione, vorrei sapere qualcosa degli altri oggetti. Potresti dirmi dove si trovano i tuoi "T", "O", "S", "D", "A" e "M"? Questo mi aiuterà a capire se la posizione centrale che hai proposto è fattibile anche per gli altri oggetti. 
	  } 
	   } 
	   } 
	  \\ 
 

    \theutterance \stepcounter{utterance}  

    & & & \multicolumn{2}{p{0.3\linewidth}}{\cellcolor[rgb]{0.95,0.95,0.95}{%
	\makecell[{{p{\linewidth}}}]{% 
	  \tt {\tiny [GM$|$GM]}  
	 COMUNICA: Sono d'accordo, spostare i "P" in una posizione centrale potrebbe essere un buon inizio. Tuttavia, prima di prendere una decisione, vorrei sapere qualcosa degli altri oggetti. Potresti dirmi dove si trovano i tuoi "T", "O", "S", "D", "A" e "M"? Questo mi aiuterà a capire se la posizione centrale che hai proposto è fattibile anche per gli altri oggetti. 
	  } 
	   } 
	   } 
	 & & \\ 
 

    \theutterance \stepcounter{utterance}  

    & & \multicolumn{4}{p{0.6\linewidth}}{\cellcolor[rgb]{0.9,0.9,0.9}{%
	\makecell[{{p{\linewidth}}}]{% 
	  \tt {\tiny [A$\langle$GM]}  
	 Il tuo messaggio è stato inoltrato all'altro giocatore.\\ \tt \\ \tt \\ \tt You are currently playing round 2 of maximum 28.\\ \tt You have collectively accumulated 0 of 16 penalties.\\ \tt L'altro giocatore ha inviato questo messaggio:\\ \tt \\ \tt \\ \tt "Sono d'accordo, spostare i "P" in una posizione centrale potrebbe essere un buon inizio. Tuttavia, prima di prendere una decisione, vorrei sapere qualcosa degli altri oggetti. Potresti dirmi dove si trovano i tuoi "T", "O", "S", "D", "A" e "M"? Questo mi aiuterà a capire se la posizione centrale che hai proposto è fattibile anche per gli altri oggetti."\\ \tt \\ \tt \\ \tt Qual è il tuo prossimo comando? 
	  } 
	   } 
	   } 
	 & \\ 
 

    \theutterance \stepcounter{utterance}  

    & \multicolumn{4}{p{0.6\linewidth}}{\cellcolor[rgb]{0.8,1,0.9}{%
	\makecell[{{p{\linewidth}}}]{% 
	  \tt {\tiny [A$\rangle$GM]}  
	 COMUNICA: Okay, cercherò di descrivere la disposizione dei miei oggetti. Il mio "T" si trova alla posizione (7,7), il mio "O" è alla posizione (5,5), il mio "S" è alla posizione (3,3), il mio "D" è alla posizione (6,6), il mio "A" è alla posizione (3,3) accanto al mio "S", e il mio "M" è alla posizione (6,6) accanto al mio "D". Sembra che alcuni dei miei oggetti siano già vicini tra loro, il che potrebbe essere utile. Ora che ti ho detto la disposizione dei miei oggetti, potresti dirmi qualcosa della disposizione dei tuoi oggetti? In particolare, dove si trovano il tuo "T", "O", "S", "D", "A" e "M" rispetto al tuo "P" in (4,1)? 
	  } 
	   } 
	   } 
	 & & \\ 
 

    \theutterance \stepcounter{utterance}  

    & & & \multicolumn{2}{p{0.3\linewidth}}{\cellcolor[rgb]{0.95,0.95,0.95}{%
	\makecell[{{p{\linewidth}}}]{% 
	  \tt {\tiny [GM$|$GM]}  
	 COMUNICA: Okay, cercherò di descrivere la disposizione dei miei oggetti. Il mio "T" si trova alla posizione (7,7), il mio "O" è alla posizione (5,5), il mio "S" è alla posizione (3,3), il mio "D" è alla posizione (6,6), il mio "A" è alla posizione (3,3) accanto al mio "S", e il mio "M" è alla posizione (6,6) accanto al mio "D". Sembra che alcuni dei miei oggetti siano già vicini tra loro, il che potrebbe essere utile. Ora che ti ho detto la disposizione dei miei oggetti, potresti dirmi qualcosa della disposizione dei tuoi oggetti? In particolare, dove si trovano il tuo "T", "O", "S", "D", "A" e "M" rispetto al tuo "P" in (4,1)? 
	  } 
	   } 
	   } 
	 & & \\ 
 

    \theutterance \stepcounter{utterance}  

    & & \multicolumn{4}{p{0.6\linewidth}}{\cellcolor[rgb]{0.9,0.9,0.9}{%
	\makecell[{{p{\linewidth}}}]{% 
	  \tt {\tiny [GM$\rangle$B]}  
	 Il tuo messaggio è stato inoltrato all'altro giocatore.\\ \tt \\ \tt \\ \tt You are currently playing round 3 of maximum 28.\\ \tt You have collectively accumulated 0 of 16 penalties.\\ \tt L'altro giocatore ha inviato questo messaggio:\\ \tt \\ \tt \\ \tt "Okay, cercherò di descrivere la disposizione dei miei oggetti. Il mio "T" si trova alla posizione (7,7), il mio "O" è alla posizione (5,5), il mio "S" è alla posizione (3,3), il mio "D" è alla posizione (6,6), il mio "A" è alla posizione (3,3) accanto al mio "S", e il mio "M" è alla posizione (6,6) accanto al mio "D". Sembra che alcuni dei miei oggetti siano già vicini tra loro, il che potrebbe essere utile. Ora che ti ho detto la disposizione dei miei oggetti, potresti dirmi qualcosa della disposizione dei tuoi oggetti? In particolare, dove si trovano il tuo "T", "O", "S", "D", "A" e "M" rispetto al tuo "P" in (4,1)?"\\ \tt \\ \tt \\ \tt Qual è il tuo prossimo comando? 
	  } 
	   } 
	   } 
	 & \\ 
 

    \theutterance \stepcounter{utterance}  

    & & & \multicolumn{4}{p{0.6\linewidth}}{\cellcolor[rgb]{1,0.85,0.72}{%
	\makecell[{{p{\linewidth}}}]{% 
	  \tt {\tiny [GM$\langle$B]}  
	 COMUNICA: Grazie per la descrizione della disposizione dei tuoi oggetti! Ora posso darti una idea della disposizione dei miei oggetti. Il mio "T" si trova alla posizione (2,2), il mio "O" è alla posizione (5,5), il mio "S" è alla posizione (6,6), il mio "D" è alla posizione (2,5), il mio "A" è alla posizione (1,7) e il mio "M" è alla posizione (3,3). Nota che alcuni dei miei oggetti sono già vicini tra loro, come il mio "P" in (4,1) e il mio "M" in (3,3). \\ \tt \\ \tt Considerando le posizioni che hai descritto per i tuoi oggetti e le posizioni dei miei oggetti, potremmo iniziare a pensare a una strategia per allineare gli oggetti identici. Sembra che il tuo "S" e il mio "S" siano in posizioni diverse, quindi potremmo dover lavorare per spostarli in una posizione comune. \\ \tt \\ \tt Qual è la tua proposta per procedere? Dovremmo iniziare a spostare gli oggetti uno per uno o hai un'idea più generale per allineare gli oggetti identici? 
	  } 
	   } 
	   } 
	  \\ 
 

    \theutterance \stepcounter{utterance}  

    & & & \multicolumn{2}{p{0.3\linewidth}}{\cellcolor[rgb]{0.95,0.95,0.95}{%
	\makecell[{{p{\linewidth}}}]{% 
	  \tt {\tiny [GM$|$GM]}  
	 COMUNICA: Grazie per la descrizione della disposizione dei tuoi oggetti! Ora posso darti una idea della disposizione dei miei oggetti. Il mio "T" si trova alla posizione (2,2), il mio "O" è alla posizione (5,5), il mio "S" è alla posizione (6,6), il mio "D" è alla posizione (2,5), il mio "A" è alla posizione (1,7) e il mio "M" è alla posizione (3,3). Nota che alcuni dei miei oggetti sono già vicini tra loro, come il mio "P" in (4,1) e il mio "M" in (3,3). \\ \tt \\ \tt Considerando le posizioni che hai descritto per i tuoi oggetti e le posizioni dei miei oggetti, potremmo iniziare a pensare a una strategia per allineare gli oggetti identici. Sembra che il tuo "S" e il mio "S" siano in posizioni diverse, quindi potremmo dover lavorare per spostarli in una posizione comune. \\ \tt \\ \tt Qual è la tua proposta per procedere? Dovremmo iniziare a spostare gli oggetti uno per uno o hai un'idea più generale per allineare gli oggetti identici? 
	  } 
	   } 
	   } 
	 & & \\ 
 

    \theutterance \stepcounter{utterance}  

    & & & \multicolumn{2}{p{0.3\linewidth}}{\cellcolor[rgb]{0.95,0.95,0.95}{%
	\makecell[{{p{\linewidth}}}]{% 
	  \tt {\tiny [GM$|$GM]}  
	 Invalid format: tail is not empty: '\\ \tt \\ \tt Considerando le posizioni che hai descritto per i tuoi oggetti e le posizioni dei miei oggetti, potremmo iniziare a pensare a una strategia per allineare gli oggetti identici. Sembra che il tuo "S" e il mio "S" siano in posizioni diverse, quindi potremmo dover lavorare per spostarli in una posizione comune. \\ \tt \\ \tt Qual è la tua proposta per procedere? Dovremmo iniziare a spostare gli oggetti uno per uno o hai un'idea più generale per allineare gli oggetti identici? 
	  } 
	   } 
	   } 
	 & & \\ 
 

    \theutterance \stepcounter{utterance}  

    & & \multicolumn{4}{p{0.6\linewidth}}{\cellcolor[rgb]{0.9,0.9,0.9}{%
	\makecell[{{p{\linewidth}}}]{% 
	  \tt {\tiny [GM$\rangle$B]}  
	 Penalità: Il tuo messaggio non deve contenere nulla dopo il comando!\\ \tt \\ \tt Assicurati che la tua risposta contenga solo `COMUNICA: <MESSAGE>` o `SPOSTA: <OBJECT>, (<X>, <Y>)`, e nient'altro!\\ \tt \\ \tt You have collectively accumulated 1 of 16 penalties.\\ \tt Riprova! 
	  } 
	   } 
	   } 
	 & \\ 
 

    \theutterance \stepcounter{utterance}  

    & & & \multicolumn{4}{p{0.6\linewidth}}{\cellcolor[rgb]{1,0.85,0.72}{%
	\makecell[{{p{\linewidth}}}]{% 
	  \tt {\tiny [GM$\langle$B]}  
	 COMUNICA: Grazie per la descrizione della disposizione dei tuoi oggetti! Il mio "T" si trova alla posizione (2,2), il mio "O" è alla posizione (5,5), il mio "S" è alla posizione (6,6), il mio "D" è alla posizione (2,5), il mio "A" è alla posizione (1,7) e il mio "M" è alla posizione (3,3). Considerando le posizioni che hai descritto per i tuoi oggetti e le posizioni dei miei oggetti, potremmo iniziare a pensare a una strategia per allineare gli oggetti identici. Sembra che il tuo "S" e il mio "S" siano in posizioni diverse, quindi potremmo dover lavorare per spostarli in una posizione comune. Qual è la tua proposta per procedere? 
	  } 
	   } 
	   } 
	  \\ 
 

    \theutterance \stepcounter{utterance}  

    & & & \multicolumn{2}{p{0.3\linewidth}}{\cellcolor[rgb]{0.95,0.95,0.95}{%
	\makecell[{{p{\linewidth}}}]{% 
	  \tt {\tiny [GM$|$GM]}  
	 COMUNICA: Grazie per la descrizione della disposizione dei tuoi oggetti! Il mio "T" si trova alla posizione (2,2), il mio "O" è alla posizione (5,5), il mio "S" è alla posizione (6,6), il mio "D" è alla posizione (2,5), il mio "A" è alla posizione (1,7) e il mio "M" è alla posizione (3,3). Considerando le posizioni che hai descritto per i tuoi oggetti e le posizioni dei miei oggetti, potremmo iniziare a pensare a una strategia per allineare gli oggetti identici. Sembra che il tuo "S" e il mio "S" siano in posizioni diverse, quindi potremmo dover lavorare per spostarli in una posizione comune. Qual è la tua proposta per procedere? 
	  } 
	   } 
	   } 
	 & & \\ 
 

    \theutterance \stepcounter{utterance}  

    & & \multicolumn{4}{p{0.6\linewidth}}{\cellcolor[rgb]{0.9,0.9,0.9}{%
	\makecell[{{p{\linewidth}}}]{% 
	  \tt {\tiny [A$\langle$GM]}  
	 Il tuo messaggio è stato inoltrato all'altro giocatore.\\ \tt \\ \tt \\ \tt You are currently playing round 3 of maximum 28.\\ \tt You have collectively accumulated 1 of 16 penalties.\\ \tt L'altro giocatore ha inviato questo messaggio:\\ \tt \\ \tt \\ \tt "Grazie per la descrizione della disposizione dei tuoi oggetti! Il mio "T" si trova alla posizione (2,2), il mio "O" è alla posizione (5,5), il mio "S" è alla posizione (6,6), il mio "D" è alla posizione (2,5), il mio "A" è alla posizione (1,7) e il mio "M" è alla posizione (3,3). Considerando le posizioni che hai descritto per i tuoi oggetti e le posizioni dei miei oggetti, potremmo iniziare a pensare a una strategia per allineare gli oggetti identici. Sembra che il tuo "S" e il mio "S" siano in posizioni diverse, quindi potremmo dover lavorare per spostarli in una posizione comune. Qual è la tua proposta per procedere?"\\ \tt \\ \tt \\ \tt Qual è il tuo prossimo comando? 
	  } 
	   } 
	   } 
	 & \\ 
 

    \theutterance \stepcounter{utterance}  

    & \multicolumn{4}{p{0.6\linewidth}}{\cellcolor[rgb]{0.8,1,0.9}{%
	\makecell[{{p{\linewidth}}}]{% 
	  \tt {\tiny [A$\rangle$GM]}  
	 COMUNICA: Grazie per la descrizione della disposizione dei tuoi oggetti! Sì, hai ragione, il mio "S" e il tuo "S" sono in posizioni diverse. Considerando che il mio "S" è già vicino al centro della griglia, alla posizione (3,3), e che il tuo "M" è anche alla posizione (3,3), potremmo provare a spostare il tuo "S" verso il centro per allinearlo con il mio "S". Inoltre, noto che il tuo "A" è alla posizione (1,7) e il mio "A" è alla posizione (3,3) accanto al mio "S", quindi potremmo anche provare a spostare il tuo "A" verso il centro per allinearlo con il mio "A". Che ne pensi? Dovremmo provare a spostare il tuo "S" e il tuo "A" verso il centro della griglia? 
	  } 
	   } 
	   } 
	 & & \\ 
 

    \theutterance \stepcounter{utterance}  

    & & & \multicolumn{2}{p{0.3\linewidth}}{\cellcolor[rgb]{0.95,0.95,0.95}{%
	\makecell[{{p{\linewidth}}}]{% 
	  \tt {\tiny [GM$|$GM]}  
	 COMUNICA: Grazie per la descrizione della disposizione dei tuoi oggetti! Sì, hai ragione, il mio "S" e il tuo "S" sono in posizioni diverse. Considerando che il mio "S" è già vicino al centro della griglia, alla posizione (3,3), e che il tuo "M" è anche alla posizione (3,3), potremmo provare a spostare il tuo "S" verso il centro per allinearlo con il mio "S". Inoltre, noto che il tuo "A" è alla posizione (1,7) e il mio "A" è alla posizione (3,3) accanto al mio "S", quindi potremmo anche provare a spostare il tuo "A" verso il centro per allinearlo con il mio "A". Che ne pensi? Dovremmo provare a spostare il tuo "S" e il tuo "A" verso il centro della griglia? 
	  } 
	   } 
	   } 
	 & & \\ 
 

    \theutterance \stepcounter{utterance}  

    & & \multicolumn{4}{p{0.6\linewidth}}{\cellcolor[rgb]{0.9,0.9,0.9}{%
	\makecell[{{p{\linewidth}}}]{% 
	  \tt {\tiny [GM$\rangle$B]}  
	 Il tuo messaggio è stato inoltrato all'altro giocatore.\\ \tt \\ \tt \\ \tt You are currently playing round 4 of maximum 28.\\ \tt You have collectively accumulated 1 of 16 penalties.\\ \tt L'altro giocatore ha inviato questo messaggio:\\ \tt \\ \tt \\ \tt "Grazie per la descrizione della disposizione dei tuoi oggetti! Sì, hai ragione, il mio "S" e il tuo "S" sono in posizioni diverse. Considerando che il mio "S" è già vicino al centro della griglia, alla posizione (3,3), e che il tuo "M" è anche alla posizione (3,3), potremmo provare a spostare il tuo "S" verso il centro per allinearlo con il mio "S". Inoltre, noto che il tuo "A" è alla posizione (1,7) e il mio "A" è alla posizione (3,3) accanto al mio "S", quindi potremmo anche provare a spostare il tuo "A" verso il centro per allinearlo con il mio "A". Che ne pensi? Dovremmo provare a spostare il tuo "S" e il tuo "A" verso il centro della griglia?"\\ \tt \\ \tt \\ \tt Qual è il tuo prossimo comando? 
	  } 
	   } 
	   } 
	 & \\ 
 

    \theutterance \stepcounter{utterance}  

    & & & \multicolumn{4}{p{0.6\linewidth}}{\cellcolor[rgb]{1,0.85,0.72}{%
	\makecell[{{p{\linewidth}}}]{% 
	  \tt {\tiny [GM$\langle$B]}  
	 COMUNICA: Sì, provare a spostare il mio "S" e il mio "A" verso il centro della griglia potrebbe essere una buona strategia. Tuttavia, prima di farlo, vorrei assicurarmi che ci sia abbastanza spazio libero intorno alla posizione (3,3) per accogliere sia il mio "S" che il mio "A". Potresti confermarmi che la posizione (3,3) e le posizioni adiacenti siano libere sul tuo tavolo da gioco? 
	  } 
	   } 
	   } 
	  \\ 
 

    \theutterance \stepcounter{utterance}  

    & & & \multicolumn{2}{p{0.3\linewidth}}{\cellcolor[rgb]{0.95,0.95,0.95}{%
	\makecell[{{p{\linewidth}}}]{% 
	  \tt {\tiny [GM$|$GM]}  
	 COMUNICA: Sì, provare a spostare il mio "S" e il mio "A" verso il centro della griglia potrebbe essere una buona strategia. Tuttavia, prima di farlo, vorrei assicurarmi che ci sia abbastanza spazio libero intorno alla posizione (3,3) per accogliere sia il mio "S" che il mio "A". Potresti confermarmi che la posizione (3,3) e le posizioni adiacenti siano libere sul tuo tavolo da gioco? 
	  } 
	   } 
	   } 
	 & & \\ 
 

    \theutterance \stepcounter{utterance}  

    & & \multicolumn{4}{p{0.6\linewidth}}{\cellcolor[rgb]{0.9,0.9,0.9}{%
	\makecell[{{p{\linewidth}}}]{% 
	  \tt {\tiny [A$\langle$GM]}  
	 Il tuo messaggio è stato inoltrato all'altro giocatore.\\ \tt \\ \tt \\ \tt You are currently playing round 4 of maximum 28.\\ \tt You have collectively accumulated 1 of 16 penalties.\\ \tt L'altro giocatore ha inviato questo messaggio:\\ \tt \\ \tt \\ \tt "Sì, provare a spostare il mio "S" e il mio "A" verso il centro della griglia potrebbe essere una buona strategia. Tuttavia, prima di farlo, vorrei assicurarmi che ci sia abbastanza spazio libero intorno alla posizione (3,3) per accogliere sia il mio "S" che il mio "A". Potresti confermarmi che la posizione (3,3) e le posizioni adiacenti siano libere sul tuo tavolo da gioco?"\\ \tt \\ \tt \\ \tt Qual è il tuo prossimo comando? 
	  } 
	   } 
	   } 
	 & \\ 
 

    \theutterance \stepcounter{utterance}  

    & \multicolumn{4}{p{0.6\linewidth}}{\cellcolor[rgb]{0.8,1,0.9}{%
	\makecell[{{p{\linewidth}}}]{% 
	  \tt {\tiny [A$\rangle$GM]}  
	 COMUNICA: Sì, posso confermarti che la posizione (3,3) sul mio tavolo da gioco è occupata dal mio "S" e dal mio "A", quindi non è libera. Tuttavia, le posizioni adiacenti (2,3), (3,2), (3,4) e (4,3) sono libere. Se spostassimo il tuo "S" e il tuo "A" verso il centro, potremmo provare a posizionarli in una di queste posizioni adiacenti. In alternativa, potremmo anche provare a spostare il mio "S" e il mio "A" per fare spazio per il tuo "S" e il tuo "A" alla posizione (3,3). Qual è la tua preferenza? Vuoi provare a spostare i tuoi oggetti verso le posizioni adiacenti o vuoi che io sposti i miei oggetti per fare spazio? 
	  } 
	   } 
	   } 
	 & & \\ 
 

    \theutterance \stepcounter{utterance}  

    & & & \multicolumn{2}{p{0.3\linewidth}}{\cellcolor[rgb]{0.95,0.95,0.95}{%
	\makecell[{{p{\linewidth}}}]{% 
	  \tt {\tiny [GM$|$GM]}  
	 COMUNICA: Sì, posso confermarti che la posizione (3,3) sul mio tavolo da gioco è occupata dal mio "S" e dal mio "A", quindi non è libera. Tuttavia, le posizioni adiacenti (2,3), (3,2), (3,4) e (4,3) sono libere. Se spostassimo il tuo "S" e il tuo "A" verso il centro, potremmo provare a posizionarli in una di queste posizioni adiacenti. In alternativa, potremmo anche provare a spostare il mio "S" e il mio "A" per fare spazio per il tuo "S" e il tuo "A" alla posizione (3,3). Qual è la tua preferenza? Vuoi provare a spostare i tuoi oggetti verso le posizioni adiacenti o vuoi che io sposti i miei oggetti per fare spazio? 
	  } 
	   } 
	   } 
	 & & \\ 
 

    \theutterance \stepcounter{utterance}  

    & & \multicolumn{4}{p{0.6\linewidth}}{\cellcolor[rgb]{0.9,0.9,0.9}{%
	\makecell[{{p{\linewidth}}}]{% 
	  \tt {\tiny [GM$\rangle$B]}  
	 Il tuo messaggio è stato inoltrato all'altro giocatore.\\ \tt \\ \tt \\ \tt You are currently playing round 5 of maximum 28.\\ \tt You have collectively accumulated 1 of 16 penalties.\\ \tt L'altro giocatore ha inviato questo messaggio:\\ \tt \\ \tt \\ \tt "Sì, posso confermarti che la posizione (3,3) sul mio tavolo da gioco è occupata dal mio "S" e dal mio "A", quindi non è libera. Tuttavia, le posizioni adiacenti (2,3), (3,2), (3,4) e (4,3) sono libere. Se spostassimo il tuo "S" e il tuo "A" verso il centro, potremmo provare a posizionarli in una di queste posizioni adiacenti. In alternativa, potremmo anche provare a spostare il mio "S" e il mio "A" per fare spazio per il tuo "S" e il tuo "A" alla posizione (3,3). Qual è la tua preferenza? Vuoi provare a spostare i tuoi oggetti verso le posizioni adiacenti o vuoi che io sposti i miei oggetti per fare spazio?"\\ \tt \\ \tt \\ \tt Qual è il tuo prossimo comando? 
	  } 
	   } 
	   } 
	 & \\ 
 

    \theutterance \stepcounter{utterance}  

    & & & \multicolumn{4}{p{0.6\linewidth}}{\cellcolor[rgb]{1,0.85,0.72}{%
	\makecell[{{p{\linewidth}}}]{% 
	  \tt {\tiny [GM$\langle$B]}  
	 COMUNICA: Preferisco provare a spostare il mio "S" nella posizione (4,3), in modo da allinearli con il tuo "S" e "A" alla posizione (3,3) e nelle posizioni adiacenti. In questo modo, potremmo iniziare a creare un gruppo di oggetti identici vicini tra loro. Che ne pensi? È una mossa fattibile sul tuo tavolo da gioco? 
	  } 
	   } 
	   } 
	  \\ 
 

    \theutterance \stepcounter{utterance}  

    & & & \multicolumn{2}{p{0.3\linewidth}}{\cellcolor[rgb]{0.95,0.95,0.95}{%
	\makecell[{{p{\linewidth}}}]{% 
	  \tt {\tiny [GM$|$GM]}  
	 COMUNICA: Preferisco provare a spostare il mio "S" nella posizione (4,3), in modo da allinearli con il tuo "S" e "A" alla posizione (3,3) e nelle posizioni adiacenti. In questo modo, potremmo iniziare a creare un gruppo di oggetti identici vicini tra loro. Che ne pensi? È una mossa fattibile sul tuo tavolo da gioco? 
	  } 
	   } 
	   } 
	 & & \\ 
 

    \theutterance \stepcounter{utterance}  

    & & \multicolumn{4}{p{0.6\linewidth}}{\cellcolor[rgb]{0.9,0.9,0.9}{%
	\makecell[{{p{\linewidth}}}]{% 
	  \tt {\tiny [A$\langle$GM]}  
	 Il tuo messaggio è stato inoltrato all'altro giocatore.\\ \tt \\ \tt \\ \tt You are currently playing round 5 of maximum 28.\\ \tt You have collectively accumulated 1 of 16 penalties.\\ \tt L'altro giocatore ha inviato questo messaggio:\\ \tt \\ \tt \\ \tt "Preferisco provare a spostare il mio "S" nella posizione (4,3), in modo da allinearli con il tuo "S" e "A" alla posizione (3,3) e nelle posizioni adiacenti. In questo modo, potremmo iniziare a creare un gruppo di oggetti identici vicini tra loro. Che ne pensi? È una mossa fattibile sul tuo tavolo da gioco?"\\ \tt \\ \tt \\ \tt Qual è il tuo prossimo comando? 
	  } 
	   } 
	   } 
	 & \\ 
 

    \theutterance \stepcounter{utterance}  

    & \multicolumn{4}{p{0.6\linewidth}}{\cellcolor[rgb]{0.8,1,0.9}{%
	\makecell[{{p{\linewidth}}}]{% 
	  \tt {\tiny [A$\rangle$GM]}  
	 COMUNICA: Sì, è una mossa fattibile sul mio tavolo da gioco! La posizione (4,3) è libera, quindi potresti spostare il tuo "S" lì senza problemi. Se lo fai, il mio "S" e il mio "A" saranno già nella posizione (3,3), e il tuo "S" sarà nella posizione (4,3), quindi avremmo un buon inizio per creare un gruppo di oggetti identici vicini tra loro. Dopo questo spostamento, potremmo valutare la disposizione degli altri oggetti e decidere la prossima mossa. Sei pronto a spostare il tuo "S" nella posizione (4,3)? 
	  } 
	   } 
	   } 
	 & & \\ 
 

    \theutterance \stepcounter{utterance}  

    & & & \multicolumn{2}{p{0.3\linewidth}}{\cellcolor[rgb]{0.95,0.95,0.95}{%
	\makecell[{{p{\linewidth}}}]{% 
	  \tt {\tiny [GM$|$GM]}  
	 COMUNICA: Sì, è una mossa fattibile sul mio tavolo da gioco! La posizione (4,3) è libera, quindi potresti spostare il tuo "S" lì senza problemi. Se lo fai, il mio "S" e il mio "A" saranno già nella posizione (3,3), e il tuo "S" sarà nella posizione (4,3), quindi avremmo un buon inizio per creare un gruppo di oggetti identici vicini tra loro. Dopo questo spostamento, potremmo valutare la disposizione degli altri oggetti e decidere la prossima mossa. Sei pronto a spostare il tuo "S" nella posizione (4,3)? 
	  } 
	   } 
	   } 
	 & & \\ 
 

    \theutterance \stepcounter{utterance}  

    & & \multicolumn{4}{p{0.6\linewidth}}{\cellcolor[rgb]{0.9,0.9,0.9}{%
	\makecell[{{p{\linewidth}}}]{% 
	  \tt {\tiny [GM$\rangle$B]}  
	 Il tuo messaggio è stato inoltrato all'altro giocatore.\\ \tt \\ \tt \\ \tt You are currently playing round 6 of maximum 28.\\ \tt You have collectively accumulated 1 of 16 penalties.\\ \tt L'altro giocatore ha inviato questo messaggio:\\ \tt \\ \tt \\ \tt "Sì, è una mossa fattibile sul mio tavolo da gioco! La posizione (4,3) è libera, quindi potresti spostare il tuo "S" lì senza problemi. Se lo fai, il mio "S" e il mio "A" saranno già nella posizione (3,3), e il tuo "S" sarà nella posizione (4,3), quindi avremmo un buon inizio per creare un gruppo di oggetti identici vicini tra loro. Dopo questo spostamento, potremmo valutare la disposizione degli altri oggetti e decidere la prossima mossa. Sei pronto a spostare il tuo "S" nella posizione (4,3)?"\\ \tt \\ \tt \\ \tt Qual è il tuo prossimo comando? 
	  } 
	   } 
	   } 
	 & \\ 
 

    \theutterance \stepcounter{utterance}  

    & & & \multicolumn{4}{p{0.6\linewidth}}{\cellcolor[rgb]{1,0.85,0.72}{%
	\makecell[{{p{\linewidth}}}]{% 
	  \tt {\tiny [GM$\langle$B]}  
	 SPOSTA: S, (4,3) 
	  } 
	   } 
	   } 
	  \\ 
 

    \theutterance \stepcounter{utterance}  

    & & & \multicolumn{2}{p{0.3\linewidth}}{\cellcolor[rgb]{0.95,0.95,0.95}{%
	\makecell[{{p{\linewidth}}}]{% 
	  \tt {\tiny [GM$|$GM]}  
	 SPOSTA: S, (4,3) 
	  } 
	   } 
	   } 
	 & & \\ 
 

    \theutterance \stepcounter{utterance}  

    & & & \multicolumn{2}{p{0.3\linewidth}}{\cellcolor[rgb]{0.95,0.95,0.95}{%
	\makecell[{{p{\linewidth}}}]{% 
	  \tt {\tiny [GM$|$GM]}  
	 Penalità: (4,3) non è vuoto, ma contiene '│'.\\ \tt \\ \tt You have collectively accumulated 2 of 16 penalties.\\ \tt Riprova! 
	  } 
	   } 
	   } 
	 & & \\ 
 

    \theutterance \stepcounter{utterance}  

    & & \multicolumn{4}{p{0.6\linewidth}}{\cellcolor[rgb]{0.9,0.9,0.9}{%
	\makecell[{{p{\linewidth}}}]{% 
	  \tt {\tiny [GM$\rangle$B]}  
	 Penalità: (4,3) non è vuoto, ma contiene '│'.\\ \tt \\ \tt You have collectively accumulated 2 of 16 penalties.\\ \tt Riprova! 
	  } 
	   } 
	   } 
	 & \\ 
 

    \theutterance \stepcounter{utterance}  

    & & & \multicolumn{4}{p{0.6\linewidth}}{\cellcolor[rgb]{1,0.85,0.72}{%
	\makecell[{{p{\linewidth}}}]{% 
	  \tt {\tiny [GM$\langle$B]}  
	 SPOSTA: S, (3,4) 
	  } 
	   } 
	   } 
	  \\ 
 

    \theutterance \stepcounter{utterance}  

    & & & \multicolumn{2}{p{0.3\linewidth}}{\cellcolor[rgb]{0.95,0.95,0.95}{%
	\makecell[{{p{\linewidth}}}]{% 
	  \tt {\tiny [GM$|$GM]}  
	 SPOSTA: S, (3,4) 
	  } 
	   } 
	   } 
	 & & \\ 
 

    \theutterance \stepcounter{utterance}  

    & & & \multicolumn{2}{p{0.3\linewidth}}{\cellcolor[rgb]{0.95,0.95,0.95}{%
	\makecell[{{p{\linewidth}}}]{% 
	  \tt {\tiny [GM$|$GM]}  
	 Penalità: (3,4) non è vuoto, ma contiene '┬'.\\ \tt \\ \tt You have collectively accumulated 3 of 16 penalties.\\ \tt Riprova! 
	  } 
	   } 
	   } 
	 & & \\ 
 

    \theutterance \stepcounter{utterance}  

    & & \multicolumn{4}{p{0.6\linewidth}}{\cellcolor[rgb]{0.9,0.9,0.9}{%
	\makecell[{{p{\linewidth}}}]{% 
	  \tt {\tiny [GM$\rangle$B]}  
	 Penalità: (3,4) non è vuoto, ma contiene '┬'.\\ \tt \\ \tt You have collectively accumulated 3 of 16 penalties.\\ \tt Riprova! 
	  } 
	   } 
	   } 
	 & \\ 
 

    \theutterance \stepcounter{utterance}  

    & & & \multicolumn{4}{p{0.6\linewidth}}{\cellcolor[rgb]{1,0.85,0.72}{%
	\makecell[{{p{\linewidth}}}]{% 
	  \tt {\tiny [GM$\langle$B]}  
	 SPOSTA: S, (4,4) 
	  } 
	   } 
	   } 
	  \\ 
 

    \theutterance \stepcounter{utterance}  

    & & & \multicolumn{2}{p{0.3\linewidth}}{\cellcolor[rgb]{0.95,0.95,0.95}{%
	\makecell[{{p{\linewidth}}}]{% 
	  \tt {\tiny [GM$|$GM]}  
	 SPOSTA: S, (4,4) 
	  } 
	   } 
	   } 
	 & & \\ 
 

    \theutterance \stepcounter{utterance}  

    & & & \multicolumn{2}{p{0.3\linewidth}}{\cellcolor[rgb]{0.95,0.95,0.95}{%
	\makecell[{{p{\linewidth}}}]{% 
	  \tt {\tiny [GM$|$GM]}  
	 Penalità: (4,4) non è vuoto, ma contiene '┴'.\\ \tt \\ \tt You have collectively accumulated 4 of 16 penalties.\\ \tt Riprova! 
	  } 
	   } 
	   } 
	 & & \\ 
 

    \theutterance \stepcounter{utterance}  

    & & \multicolumn{4}{p{0.6\linewidth}}{\cellcolor[rgb]{0.9,0.9,0.9}{%
	\makecell[{{p{\linewidth}}}]{% 
	  \tt {\tiny [GM$\rangle$B]}  
	 Penalità: (4,4) non è vuoto, ma contiene '┴'.\\ \tt \\ \tt You have collectively accumulated 4 of 16 penalties.\\ \tt Riprova! 
	  } 
	   } 
	   } 
	 & \\ 
 

    \theutterance \stepcounter{utterance}  

    & & & \multicolumn{4}{p{0.6\linewidth}}{\cellcolor[rgb]{1,0.85,0.72}{%
	\makecell[{{p{\linewidth}}}]{% 
	  \tt {\tiny [GM$\langle$B]}  
	 SPOSTA: S, (5,6) 
	  } 
	   } 
	   } 
	  \\ 
 

    \theutterance \stepcounter{utterance}  

    & & & \multicolumn{2}{p{0.3\linewidth}}{\cellcolor[rgb]{0.95,0.95,0.95}{%
	\makecell[{{p{\linewidth}}}]{% 
	  \tt {\tiny [GM$|$GM]}  
	 SPOSTA: S, (5,6) 
	  } 
	   } 
	   } 
	 & & \\ 
 

    \theutterance \stepcounter{utterance}  

    & & & \multicolumn{2}{p{0.3\linewidth}}{\cellcolor[rgb]{0.95,0.95,0.95}{%
	\makecell[{{p{\linewidth}}}]{% 
	  \tt {\tiny [GM$|$GM]}  
	 Penalità: (5,6) non è vuoto, ma contiene 'S'.\\ \tt \\ \tt You have collectively accumulated 5 of 16 penalties.\\ \tt Riprova! 
	  } 
	   } 
	   } 
	 & & \\ 
 

    \theutterance \stepcounter{utterance}  

    & & \multicolumn{4}{p{0.6\linewidth}}{\cellcolor[rgb]{0.9,0.9,0.9}{%
	\makecell[{{p{\linewidth}}}]{% 
	  \tt {\tiny [GM$\rangle$B]}  
	 Penalità: (5,6) non è vuoto, ma contiene 'S'.\\ \tt \\ \tt You have collectively accumulated 5 of 16 penalties.\\ \tt Riprova! 
	  } 
	   } 
	   } 
	 & \\ 
 

    \theutterance \stepcounter{utterance}  

    & & & \multicolumn{4}{p{0.6\linewidth}}{\cellcolor[rgb]{1,0.85,0.72}{%
	\makecell[{{p{\linewidth}}}]{% 
	  \tt {\tiny [GM$\langle$B]}  
	 SPOSTA: S, (5,4) 
	  } 
	   } 
	   } 
	  \\ 
 

    \theutterance \stepcounter{utterance}  

    & & & \multicolumn{2}{p{0.3\linewidth}}{\cellcolor[rgb]{0.95,0.95,0.95}{%
	\makecell[{{p{\linewidth}}}]{% 
	  \tt {\tiny [GM$|$GM]}  
	 SPOSTA: S, (5,4) 
	  } 
	   } 
	   } 
	 & & \\ 
 

    \theutterance \stepcounter{utterance}  

    & & & \multicolumn{2}{p{0.3\linewidth}}{\cellcolor[rgb]{0.95,0.95,0.95}{%
	\makecell[{{p{\linewidth}}}]{% 
	  \tt {\tiny [GM$|$GM]}  
	 Penalità: (5,4) non è vuoto, ma contiene '─'.\\ \tt \\ \tt You have collectively accumulated 6 of 16 penalties.\\ \tt Riprova! 
	  } 
	   } 
	   } 
	 & & \\ 
 

    \theutterance \stepcounter{utterance}  

    & & \multicolumn{4}{p{0.6\linewidth}}{\cellcolor[rgb]{0.9,0.9,0.9}{%
	\makecell[{{p{\linewidth}}}]{% 
	  \tt {\tiny [GM$\rangle$B]}  
	 Penalità: (5,4) non è vuoto, ma contiene '─'.\\ \tt \\ \tt You have collectively accumulated 6 of 16 penalties.\\ \tt Riprova! 
	  } 
	   } 
	   } 
	 & \\ 
 

    \theutterance \stepcounter{utterance}  

    & & & \multicolumn{4}{p{0.6\linewidth}}{\cellcolor[rgb]{1,0.85,0.72}{%
	\makecell[{{p{\linewidth}}}]{% 
	  \tt {\tiny [GM$\langle$B]}  
	 SPOSTA: S, (3,5) 
	  } 
	   } 
	   } 
	  \\ 
 

    \theutterance \stepcounter{utterance}  

    & & & \multicolumn{2}{p{0.3\linewidth}}{\cellcolor[rgb]{0.95,0.95,0.95}{%
	\makecell[{{p{\linewidth}}}]{% 
	  \tt {\tiny [GM$|$GM]}  
	 SPOSTA: S, (3,5) 
	  } 
	   } 
	   } 
	 & & \\ 
 

    \theutterance \stepcounter{utterance}  

    & & & \multicolumn{2}{p{0.3\linewidth}}{\cellcolor[rgb]{0.95,0.95,0.95}{%
	\makecell[{{p{\linewidth}}}]{% 
	  \tt {\tiny [GM$|$GM]}  
	 Penalità: (3,5) non è vuoto, ma contiene '│'.\\ \tt \\ \tt You have collectively accumulated 7 of 16 penalties.\\ \tt Riprova! 
	  } 
	   } 
	   } 
	 & & \\ 
 

    \theutterance \stepcounter{utterance}  

    & & \multicolumn{4}{p{0.6\linewidth}}{\cellcolor[rgb]{0.9,0.9,0.9}{%
	\makecell[{{p{\linewidth}}}]{% 
	  \tt {\tiny [GM$\rangle$B]}  
	 Penalità: (3,5) non è vuoto, ma contiene '│'.\\ \tt \\ \tt You have collectively accumulated 7 of 16 penalties.\\ \tt Riprova! 
	  } 
	   } 
	   } 
	 & \\ 
 

    \theutterance \stepcounter{utterance}  

    & & & \multicolumn{4}{p{0.6\linewidth}}{\cellcolor[rgb]{1,0.85,0.72}{%
	\makecell[{{p{\linewidth}}}]{% 
	  \tt {\tiny [GM$\langle$B]}  
	 SPOSTA: S, (5,3) 
	  } 
	   } 
	   } 
	  \\ 
 

    \theutterance \stepcounter{utterance}  

    & & & \multicolumn{2}{p{0.3\linewidth}}{\cellcolor[rgb]{0.95,0.95,0.95}{%
	\makecell[{{p{\linewidth}}}]{% 
	  \tt {\tiny [GM$|$GM]}  
	 SPOSTA: S, (5,3) 
	  } 
	   } 
	   } 
	 & & \\ 
 

    \theutterance \stepcounter{utterance}  

    & & & \multicolumn{2}{p{0.3\linewidth}}{\cellcolor[rgb]{0.95,0.95,0.95}{%
	\makecell[{{p{\linewidth}}}]{% 
	  \tt {\tiny [GM$|$GM]}  
	 Penalità: (5,3) non è vuoto, ma contiene 'M'.\\ \tt \\ \tt You have collectively accumulated 8 of 16 penalties.\\ \tt Riprova! 
	  } 
	   } 
	   } 
	 & & \\ 
 

    \theutterance \stepcounter{utterance}  

    & & \multicolumn{4}{p{0.6\linewidth}}{\cellcolor[rgb]{0.9,0.9,0.9}{%
	\makecell[{{p{\linewidth}}}]{% 
	  \tt {\tiny [GM$\rangle$B]}  
	 Penalità: (5,3) non è vuoto, ma contiene 'M'.\\ \tt \\ \tt You have collectively accumulated 8 of 16 penalties.\\ \tt Riprova! 
	  } 
	   } 
	   } 
	 & \\ 
 

    \theutterance \stepcounter{utterance}  

    & & & \multicolumn{4}{p{0.6\linewidth}}{\cellcolor[rgb]{1,0.85,0.72}{%
	\makecell[{{p{\linewidth}}}]{% 
	  \tt {\tiny [GM$\langle$B]}  
	 SPOSTA: S, (4,5) 
	  } 
	   } 
	   } 
	  \\ 
 

    \theutterance \stepcounter{utterance}  

    & & & \multicolumn{2}{p{0.3\linewidth}}{\cellcolor[rgb]{0.95,0.95,0.95}{%
	\makecell[{{p{\linewidth}}}]{% 
	  \tt {\tiny [GM$|$GM]}  
	 SPOSTA: S, (4,5) 
	  } 
	   } 
	   } 
	 & & \\ 
 

    \theutterance \stepcounter{utterance}  

    & & & \multicolumn{2}{p{0.3\linewidth}}{\cellcolor[rgb]{0.95,0.95,0.95}{%
	\makecell[{{p{\linewidth}}}]{% 
	  \tt {\tiny [GM$|$GM]}  
	 Spostato 'S' in (4,5) con successo. La tua griglia aggiornata appare così:\\ \tt ```\\ \tt  1234567\\ \tt ╔═╤═╤═╤═╗\\ \tt ║◌│◌│◌│P║ 1\\ \tt ║T└─┼─┴─╢ 2\\ \tt ║◌◌◌│M◌◌║ 3\\ \tt ╟──┬┴───╢ 4\\ \tt ║◌D│S◌O◌║ 5\\ \tt ║◌◌│◌◌◌◌║ 6\\ \tt ║A◌│◌◌◌◌║ 7\\ \tt ╚══╧════╝\\ \tt ```\\ \tt  
	  } 
	   } 
	   } 
	 & & \\ 
 

    \theutterance \stepcounter{utterance}  

    & & \multicolumn{4}{p{0.6\linewidth}}{\cellcolor[rgb]{0.9,0.9,0.9}{%
	\makecell[{{p{\linewidth}}}]{% 
	  \tt {\tiny [A$\langle$GM]}  
	 Il tuo messaggio è stato inoltrato all'altro giocatore.\\ \tt \\ \tt \\ \tt You are currently playing round 6 of maximum 28.\\ \tt You have collectively accumulated 8 of 16 penalties.\\ \tt L'altro giocatore ha spostato un oggetto sulla sua griglia.\\ \tt \\ \tt \\ \tt Qual è il tuo prossimo comando? 
	  } 
	   } 
	   } 
	 & \\ 
 

    \theutterance \stepcounter{utterance}  

    & \multicolumn{4}{p{0.6\linewidth}}{\cellcolor[rgb]{0.8,1,0.9}{%
	\makecell[{{p{\linewidth}}}]{% 
	  \tt {\tiny [A$\rangle$GM]}  
	 COMUNICA: Sembra che ci siano state alcune penalità accumulate. Spero che non ci siano problemi gravi. Comunque, visto che l'altro giocatore ha spostato un oggetto sulla sua griglia, vorrei sapere se il mio "S" e il tuo "S" sono ora allineati nella posizione (3,3) e (4,3) come pianificato. Potresti confermarmi se lo spostamento è stato eseguito correttamente e se ci sono altri oggetti che potrebbero essere spostati per allinearsi con quelli già posizionati? 
	  } 
	   } 
	   } 
	 & & \\ 
 

    \theutterance \stepcounter{utterance}  

    & & & \multicolumn{2}{p{0.3\linewidth}}{\cellcolor[rgb]{0.95,0.95,0.95}{%
	\makecell[{{p{\linewidth}}}]{% 
	  \tt {\tiny [GM$|$GM]}  
	 COMUNICA: Sembra che ci siano state alcune penalità accumulate. Spero che non ci siano problemi gravi. Comunque, visto che l'altro giocatore ha spostato un oggetto sulla sua griglia, vorrei sapere se il mio "S" e il tuo "S" sono ora allineati nella posizione (3,3) e (4,3) come pianificato. Potresti confermarmi se lo spostamento è stato eseguito correttamente e se ci sono altri oggetti che potrebbero essere spostati per allinearsi con quelli già posizionati? 
	  } 
	   } 
	   } 
	 & & \\ 
 

    \theutterance \stepcounter{utterance}  

    & & \multicolumn{4}{p{0.6\linewidth}}{\cellcolor[rgb]{0.9,0.9,0.9}{%
	\makecell[{{p{\linewidth}}}]{% 
	  \tt {\tiny [GM$\rangle$B]}  
	 Spostato 'S' in (4,5) con successo. La tua griglia aggiornata appare così:\\ \tt ```\\ \tt  1234567\\ \tt ╔═╤═╤═╤═╗\\ \tt ║◌│◌│◌│P║ 1\\ \tt ║T└─┼─┴─╢ 2\\ \tt ║◌◌◌│M◌◌║ 3\\ \tt ╟──┬┴───╢ 4\\ \tt ║◌D│S◌O◌║ 5\\ \tt ║◌◌│◌◌◌◌║ 6\\ \tt ║A◌│◌◌◌◌║ 7\\ \tt ╚══╧════╝\\ \tt ```\\ \tt You are currently playing round 7 of maximum 28.\\ \tt You have collectively accumulated 8 of 16 penalties.\\ \tt L'altro giocatore ha inviato questo messaggio:\\ \tt \\ \tt \\ \tt "Sembra che ci siano state alcune penalità accumulate. Spero che non ci siano problemi gravi. Comunque, visto che l'altro giocatore ha spostato un oggetto sulla sua griglia, vorrei sapere se il mio "S" e il tuo "S" sono ora allineati nella posizione (3,3) e (4,3) come pianificato. Potresti confermarmi se lo spostamento è stato eseguito correttamente e se ci sono altri oggetti che potrebbero essere spostati per allinearsi con quelli già posizionati?"\\ \tt \\ \tt \\ \tt Qual è il tuo prossimo comando? 
	  } 
	   } 
	   } 
	 & \\ 
 

    \theutterance \stepcounter{utterance}  

    & & & \multicolumn{4}{p{0.6\linewidth}}{\cellcolor[rgb]{1,0.85,0.72}{%
	\makecell[{{p{\linewidth}}}]{% 
	  \tt {\tiny [GM$\langle$B]}  
	 COMUNICA: Sì, ho spostato il mio "S" nella posizione (4,5), ma non nella posizione (4,3) come pianificato, a causa di alcuni problemi con le posizioni occupate. Adesso che il mio "S" è nella posizione (4,5), potremmo valutare come allineare gli altri oggetti. Il mio "A" è ancora nella posizione (1,7) e il mio "D" è nella posizione (2,5). Potremmo provare a spostare alcuni di questi oggetti per allinearsi con il mio "S" e il tuo "S". Che ne pensi? Dovremmo continuare a lavorare sull'allineamento degli oggetti? 
	  } 
	   } 
	   } 
	  \\ 
 

    \theutterance \stepcounter{utterance}  

    & & & \multicolumn{2}{p{0.3\linewidth}}{\cellcolor[rgb]{0.95,0.95,0.95}{%
	\makecell[{{p{\linewidth}}}]{% 
	  \tt {\tiny [GM$|$GM]}  
	 COMUNICA: Sì, ho spostato il mio "S" nella posizione (4,5), ma non nella posizione (4,3) come pianificato, a causa di alcuni problemi con le posizioni occupate. Adesso che il mio "S" è nella posizione (4,5), potremmo valutare come allineare gli altri oggetti. Il mio "A" è ancora nella posizione (1,7) e il mio "D" è nella posizione (2,5). Potremmo provare a spostare alcuni di questi oggetti per allinearsi con il mio "S" e il tuo "S". Che ne pensi? Dovremmo continuare a lavorare sull'allineamento degli oggetti? 
	  } 
	   } 
	   } 
	 & & \\ 
 

    \theutterance \stepcounter{utterance}  

    & & \multicolumn{4}{p{0.6\linewidth}}{\cellcolor[rgb]{0.9,0.9,0.9}{%
	\makecell[{{p{\linewidth}}}]{% 
	  \tt {\tiny [A$\langle$GM]}  
	 Il tuo messaggio è stato inoltrato all'altro giocatore.\\ \tt \\ \tt \\ \tt You are currently playing round 7 of maximum 28.\\ \tt You have collectively accumulated 8 of 16 penalties.\\ \tt L'altro giocatore ha inviato questo messaggio:\\ \tt \\ \tt \\ \tt "Sì, ho spostato il mio "S" nella posizione (4,5), ma non nella posizione (4,3) come pianificato, a causa di alcuni problemi con le posizioni occupate. Adesso che il mio "S" è nella posizione (4,5), potremmo valutare come allineare gli altri oggetti. Il mio "A" è ancora nella posizione (1,7) e il mio "D" è nella posizione (2,5). Potremmo provare a spostare alcuni di questi oggetti per allinearsi con il mio "S" e il tuo "S". Che ne pensi? Dovremmo continuare a lavorare sull'allineamento degli oggetti?"\\ \tt \\ \tt \\ \tt Qual è il tuo prossimo comando? 
	  } 
	   } 
	   } 
	 & \\ 
 

    \theutterance \stepcounter{utterance}  

    & \multicolumn{4}{p{0.6\linewidth}}{\cellcolor[rgb]{0.8,1,0.9}{%
	\makecell[{{p{\linewidth}}}]{% 
	  \tt {\tiny [A$\rangle$GM]}  
	 COMUNICA: Ah, capisco che ci sono stati alcuni problemi con le posizioni occupate. Grazie per la spiegazione. Sì, sono d'accordo che continueremo a lavorare sull'allineamento degli oggetti. Considerando la nuova posizione del tuo "S" (4,5), potremmo valutare se spostare il mio "S" o altri oggetti per allinearli con esso. Inoltre, visto che il tuo "A" è nella posizione (1,7) e il mio "A" è vicino al mio "S", potremmo provare a spostare il tuo "A" verso il centro per allinearli. Il tuo "D" è nella posizione (2,5), che è vicina alla posizione del tuo "S", quindi potremmo anche valutare se spostarlo per allinearli. Che ne pensi di spostare il tuo "A" nella posizione (3,5) o (3,4) per iniziare a creare un gruppo di oggetti identici vicini tra loro? 
	  } 
	   } 
	   } 
	 & & \\ 
 

    \theutterance \stepcounter{utterance}  

    & & & \multicolumn{2}{p{0.3\linewidth}}{\cellcolor[rgb]{0.95,0.95,0.95}{%
	\makecell[{{p{\linewidth}}}]{% 
	  \tt {\tiny [GM$|$GM]}  
	 COMUNICA: Ah, capisco che ci sono stati alcuni problemi con le posizioni occupate. Grazie per la spiegazione. Sì, sono d'accordo che continueremo a lavorare sull'allineamento degli oggetti. Considerando la nuova posizione del tuo "S" (4,5), potremmo valutare se spostare il mio "S" o altri oggetti per allinearli con esso. Inoltre, visto che il tuo "A" è nella posizione (1,7) e il mio "A" è vicino al mio "S", potremmo provare a spostare il tuo "A" verso il centro per allinearli. Il tuo "D" è nella posizione (2,5), che è vicina alla posizione del tuo "S", quindi potremmo anche valutare se spostarlo per allinearli. Che ne pensi di spostare il tuo "A" nella posizione (3,5) o (3,4) per iniziare a creare un gruppo di oggetti identici vicini tra loro? 
	  } 
	   } 
	   } 
	 & & \\ 
 

    \theutterance \stepcounter{utterance}  

    & & \multicolumn{4}{p{0.6\linewidth}}{\cellcolor[rgb]{0.9,0.9,0.9}{%
	\makecell[{{p{\linewidth}}}]{% 
	  \tt {\tiny [GM$\rangle$B]}  
	 Il tuo messaggio è stato inoltrato all'altro giocatore.\\ \tt \\ \tt \\ \tt You are currently playing round 8 of maximum 28.\\ \tt You have collectively accumulated 8 of 16 penalties.\\ \tt L'altro giocatore ha inviato questo messaggio:\\ \tt \\ \tt \\ \tt "Ah, capisco che ci sono stati alcuni problemi con le posizioni occupate. Grazie per la spiegazione. Sì, sono d'accordo che continueremo a lavorare sull'allineamento degli oggetti. Considerando la nuova posizione del tuo "S" (4,5), potremmo valutare se spostare il mio "S" o altri oggetti per allinearli con esso. Inoltre, visto che il tuo "A" è nella posizione (1,7) e il mio "A" è vicino al mio "S", potremmo provare a spostare il tuo "A" verso il centro per allinearli. Il tuo "D" è nella posizione (2,5), che è vicina alla posizione del tuo "S", quindi potremmo anche valutare se spostarlo per allinearli. Che ne pensi di spostare il tuo "A" nella posizione (3,5) o (3,4) per iniziare a creare un gruppo di oggetti identici vicini tra loro?"\\ \tt \\ \tt \\ \tt Qual è il tuo prossimo comando? 
	  } 
	   } 
	   } 
	 & \\ 
 

    \theutterance \stepcounter{utterance}  

    & & & \multicolumn{4}{p{0.6\linewidth}}{\cellcolor[rgb]{1,0.85,0.72}{%
	\makecell[{{p{\linewidth}}}]{% 
	  \tt {\tiny [GM$\langle$B]}  
	 SPOSTA: A, (3,5) 
	  } 
	   } 
	   } 
	  \\ 
 

    \theutterance \stepcounter{utterance}  

    & & & \multicolumn{2}{p{0.3\linewidth}}{\cellcolor[rgb]{0.95,0.95,0.95}{%
	\makecell[{{p{\linewidth}}}]{% 
	  \tt {\tiny [GM$|$GM]}  
	 SPOSTA: A, (3,5) 
	  } 
	   } 
	   } 
	 & & \\ 
 

    \theutterance \stepcounter{utterance}  

    & & & \multicolumn{2}{p{0.3\linewidth}}{\cellcolor[rgb]{0.95,0.95,0.95}{%
	\makecell[{{p{\linewidth}}}]{% 
	  \tt {\tiny [GM$|$GM]}  
	 Penalità: (3,5) non è vuoto, ma contiene '│'.\\ \tt \\ \tt You have collectively accumulated 9 of 16 penalties.\\ \tt Riprova! 
	  } 
	   } 
	   } 
	 & & \\ 
 

    \theutterance \stepcounter{utterance}  

    & & \multicolumn{4}{p{0.6\linewidth}}{\cellcolor[rgb]{0.9,0.9,0.9}{%
	\makecell[{{p{\linewidth}}}]{% 
	  \tt {\tiny [GM$\rangle$B]}  
	 Penalità: (3,5) non è vuoto, ma contiene '│'.\\ \tt \\ \tt You have collectively accumulated 9 of 16 penalties.\\ \tt Riprova! 
	  } 
	   } 
	   } 
	 & \\ 
 

    \theutterance \stepcounter{utterance}  

    & & & \multicolumn{4}{p{0.6\linewidth}}{\cellcolor[rgb]{1,0.85,0.72}{%
	\makecell[{{p{\linewidth}}}]{% 
	  \tt {\tiny [GM$\langle$B]}  
	 SPOSTA: A, (3,6) 
	  } 
	   } 
	   } 
	  \\ 
 

    \theutterance \stepcounter{utterance}  

    & & & \multicolumn{2}{p{0.3\linewidth}}{\cellcolor[rgb]{0.95,0.95,0.95}{%
	\makecell[{{p{\linewidth}}}]{% 
	  \tt {\tiny [GM$|$GM]}  
	 SPOSTA: A, (3,6) 
	  } 
	   } 
	   } 
	 & & \\ 
 

    \theutterance \stepcounter{utterance}  

    & & & \multicolumn{2}{p{0.3\linewidth}}{\cellcolor[rgb]{0.95,0.95,0.95}{%
	\makecell[{{p{\linewidth}}}]{% 
	  \tt {\tiny [GM$|$GM]}  
	 Penalità: (3,6) non è vuoto, ma contiene '│'.\\ \tt \\ \tt You have collectively accumulated 10 of 16 penalties.\\ \tt Riprova! 
	  } 
	   } 
	   } 
	 & & \\ 
 

    \theutterance \stepcounter{utterance}  

    & & \multicolumn{4}{p{0.6\linewidth}}{\cellcolor[rgb]{0.9,0.9,0.9}{%
	\makecell[{{p{\linewidth}}}]{% 
	  \tt {\tiny [GM$\rangle$B]}  
	 Penalità: (3,6) non è vuoto, ma contiene '│'.\\ \tt \\ \tt You have collectively accumulated 10 of 16 penalties.\\ \tt Riprova! 
	  } 
	   } 
	   } 
	 & \\ 
 

    \theutterance \stepcounter{utterance}  

    & & & \multicolumn{4}{p{0.6\linewidth}}{\cellcolor[rgb]{1,0.85,0.72}{%
	\makecell[{{p{\linewidth}}}]{% 
	  \tt {\tiny [GM$\langle$B]}  
	 SPOSTA: A, (3,4) 
	  } 
	   } 
	   } 
	  \\ 
 

    \theutterance \stepcounter{utterance}  

    & & & \multicolumn{2}{p{0.3\linewidth}}{\cellcolor[rgb]{0.95,0.95,0.95}{%
	\makecell[{{p{\linewidth}}}]{% 
	  \tt {\tiny [GM$|$GM]}  
	 SPOSTA: A, (3,4) 
	  } 
	   } 
	   } 
	 & & \\ 
 

    \theutterance \stepcounter{utterance}  

    & & & \multicolumn{2}{p{0.3\linewidth}}{\cellcolor[rgb]{0.95,0.95,0.95}{%
	\makecell[{{p{\linewidth}}}]{% 
	  \tt {\tiny [GM$|$GM]}  
	 Penalità: (3,4) non è vuoto, ma contiene '┬'.\\ \tt \\ \tt You have collectively accumulated 11 of 16 penalties.\\ \tt Riprova! 
	  } 
	   } 
	   } 
	 & & \\ 
 

    \theutterance \stepcounter{utterance}  

    & & \multicolumn{4}{p{0.6\linewidth}}{\cellcolor[rgb]{0.9,0.9,0.9}{%
	\makecell[{{p{\linewidth}}}]{% 
	  \tt {\tiny [GM$\rangle$B]}  
	 Penalità: (3,4) non è vuoto, ma contiene '┬'.\\ \tt \\ \tt You have collectively accumulated 11 of 16 penalties.\\ \tt Riprova! 
	  } 
	   } 
	   } 
	 & \\ 
 

    \theutterance \stepcounter{utterance}  

    & & & \multicolumn{4}{p{0.6\linewidth}}{\cellcolor[rgb]{1,0.85,0.72}{%
	\makecell[{{p{\linewidth}}}]{% 
	  \tt {\tiny [GM$\langle$B]}  
	 SPOSTA: A, (4,4) 
	  } 
	   } 
	   } 
	  \\ 
 

    \theutterance \stepcounter{utterance}  

    & & & \multicolumn{2}{p{0.3\linewidth}}{\cellcolor[rgb]{0.95,0.95,0.95}{%
	\makecell[{{p{\linewidth}}}]{% 
	  \tt {\tiny [GM$|$GM]}  
	 SPOSTA: A, (4,4) 
	  } 
	   } 
	   } 
	 & & \\ 
 

    \theutterance \stepcounter{utterance}  

    & & & \multicolumn{2}{p{0.3\linewidth}}{\cellcolor[rgb]{0.95,0.95,0.95}{%
	\makecell[{{p{\linewidth}}}]{% 
	  \tt {\tiny [GM$|$GM]}  
	 Penalità: (4,4) non è vuoto, ma contiene '┴'.\\ \tt \\ \tt You have collectively accumulated 12 of 16 penalties.\\ \tt Riprova! 
	  } 
	   } 
	   } 
	 & & \\ 
 

    \theutterance \stepcounter{utterance}  

    & & \multicolumn{4}{p{0.6\linewidth}}{\cellcolor[rgb]{0.9,0.9,0.9}{%
	\makecell[{{p{\linewidth}}}]{% 
	  \tt {\tiny [GM$\rangle$B]}  
	 Penalità: (4,4) non è vuoto, ma contiene '┴'.\\ \tt \\ \tt You have collectively accumulated 12 of 16 penalties.\\ \tt Riprova! 
	  } 
	   } 
	   } 
	 & \\ 
 

    \theutterance \stepcounter{utterance}  

    & & & \multicolumn{4}{p{0.6\linewidth}}{\cellcolor[rgb]{1,0.85,0.72}{%
	\makecell[{{p{\linewidth}}}]{% 
	  \tt {\tiny [GM$\langle$B]}  
	 SPOSTA: A, (5,2) 
	  } 
	   } 
	   } 
	  \\ 
 

    \theutterance \stepcounter{utterance}  

    & & & \multicolumn{2}{p{0.3\linewidth}}{\cellcolor[rgb]{0.95,0.95,0.95}{%
	\makecell[{{p{\linewidth}}}]{% 
	  \tt {\tiny [GM$|$GM]}  
	 SPOSTA: A, (5,2) 
	  } 
	   } 
	   } 
	 & & \\ 
 

    \theutterance \stepcounter{utterance}  

    & & & \multicolumn{2}{p{0.3\linewidth}}{\cellcolor[rgb]{0.95,0.95,0.95}{%
	\makecell[{{p{\linewidth}}}]{% 
	  \tt {\tiny [GM$|$GM]}  
	 Penalità: (5,2) non è vuoto, ma contiene '─'.\\ \tt \\ \tt You have collectively accumulated 13 of 16 penalties.\\ \tt Riprova! 
	  } 
	   } 
	   } 
	 & & \\ 
 

    \theutterance \stepcounter{utterance}  

    & & \multicolumn{4}{p{0.6\linewidth}}{\cellcolor[rgb]{0.9,0.9,0.9}{%
	\makecell[{{p{\linewidth}}}]{% 
	  \tt {\tiny [GM$\rangle$B]}  
	 Penalità: (5,2) non è vuoto, ma contiene '─'.\\ \tt \\ \tt You have collectively accumulated 13 of 16 penalties.\\ \tt Riprova! 
	  } 
	   } 
	   } 
	 & \\ 
 

    \theutterance \stepcounter{utterance}  

    & & & \multicolumn{4}{p{0.6\linewidth}}{\cellcolor[rgb]{1,0.85,0.72}{%
	\makecell[{{p{\linewidth}}}]{% 
	  \tt {\tiny [GM$\langle$B]}  
	 SPOSTA: A, (2,4) 
	  } 
	   } 
	   } 
	  \\ 
 

    \theutterance \stepcounter{utterance}  

    & & & \multicolumn{2}{p{0.3\linewidth}}{\cellcolor[rgb]{0.95,0.95,0.95}{%
	\makecell[{{p{\linewidth}}}]{% 
	  \tt {\tiny [GM$|$GM]}  
	 SPOSTA: A, (2,4) 
	  } 
	   } 
	   } 
	 & & \\ 
 

    \theutterance \stepcounter{utterance}  

    & & & \multicolumn{2}{p{0.3\linewidth}}{\cellcolor[rgb]{0.95,0.95,0.95}{%
	\makecell[{{p{\linewidth}}}]{% 
	  \tt {\tiny [GM$|$GM]}  
	 Penalità: (2,4) non è vuoto, ma contiene '─'.\\ \tt \\ \tt You have collectively accumulated 14 of 16 penalties.\\ \tt Riprova! 
	  } 
	   } 
	   } 
	 & & \\ 
 

    \theutterance \stepcounter{utterance}  

    & & \multicolumn{4}{p{0.6\linewidth}}{\cellcolor[rgb]{0.9,0.9,0.9}{%
	\makecell[{{p{\linewidth}}}]{% 
	  \tt {\tiny [GM$\rangle$B]}  
	 Penalità: (2,4) non è vuoto, ma contiene '─'.\\ \tt \\ \tt You have collectively accumulated 14 of 16 penalties.\\ \tt Riprova! 
	  } 
	   } 
	   } 
	 & \\ 
 

    \theutterance \stepcounter{utterance}  

    & & & \multicolumn{4}{p{0.6\linewidth}}{\cellcolor[rgb]{1,0.85,0.72}{%
	\makecell[{{p{\linewidth}}}]{% 
	  \tt {\tiny [GM$\langle$B]}  
	 SPOSTA: A, (6,2) 
	  } 
	   } 
	   } 
	  \\ 
 

    \theutterance \stepcounter{utterance}  

    & & & \multicolumn{2}{p{0.3\linewidth}}{\cellcolor[rgb]{0.95,0.95,0.95}{%
	\makecell[{{p{\linewidth}}}]{% 
	  \tt {\tiny [GM$|$GM]}  
	 SPOSTA: A, (6,2) 
	  } 
	   } 
	   } 
	 & & \\ 
 

    \theutterance \stepcounter{utterance}  

    & & & \multicolumn{2}{p{0.3\linewidth}}{\cellcolor[rgb]{0.95,0.95,0.95}{%
	\makecell[{{p{\linewidth}}}]{% 
	  \tt {\tiny [GM$|$GM]}  
	 Penalità: (6,2) non è vuoto, ma contiene '┴'.\\ \tt \\ \tt You have collectively accumulated 15 of 16 penalties.\\ \tt Riprova! 
	  } 
	   } 
	   } 
	 & & \\ 
 

    \theutterance \stepcounter{utterance}  

    & & \multicolumn{4}{p{0.6\linewidth}}{\cellcolor[rgb]{0.9,0.9,0.9}{%
	\makecell[{{p{\linewidth}}}]{% 
	  \tt {\tiny [GM$\rangle$B]}  
	 Penalità: (6,2) non è vuoto, ma contiene '┴'.\\ \tt \\ \tt You have collectively accumulated 15 of 16 penalties.\\ \tt Riprova! 
	  } 
	   } 
	   } 
	 & \\ 
 

    \theutterance \stepcounter{utterance}  

    & & & \multicolumn{4}{p{0.6\linewidth}}{\cellcolor[rgb]{1,0.85,0.72}{%
	\makecell[{{p{\linewidth}}}]{% 
	  \tt {\tiny [GM$\langle$B]}  
	 COMUNICA: Sembra che stiamo avendo problemi a trovare una posizione libera per spostare il mio "A". Potremmo valutare di spostare altri oggetti per liberare spazio o trovare un'altra strategia per allineare gli oggetti identici. Che ne pensi? Dovremmo continuare a cercare una posizione per il mio "A" o provare un approccio diverso? 
	  } 
	   } 
	   } 
	  \\ 
 

    \theutterance \stepcounter{utterance}  

    & & & \multicolumn{2}{p{0.3\linewidth}}{\cellcolor[rgb]{0.95,0.95,0.95}{%
	\makecell[{{p{\linewidth}}}]{% 
	  \tt {\tiny [GM$|$GM]}  
	 COMUNICA: Sembra che stiamo avendo problemi a trovare una posizione libera per spostare il mio "A". Potremmo valutare di spostare altri oggetti per liberare spazio o trovare un'altra strategia per allineare gli oggetti identici. Che ne pensi? Dovremmo continuare a cercare una posizione per il mio "A" o provare un approccio diverso? 
	  } 
	   } 
	   } 
	 & & \\ 
 

    \theutterance \stepcounter{utterance}  

    & & \multicolumn{4}{p{0.6\linewidth}}{\cellcolor[rgb]{0.9,0.9,0.9}{%
	\makecell[{{p{\linewidth}}}]{% 
	  \tt {\tiny [A$\langle$GM]}  
	 Il tuo messaggio è stato inoltrato all'altro giocatore.\\ \tt \\ \tt \\ \tt You are currently playing round 8 of maximum 28.\\ \tt You have collectively accumulated 15 of 16 penalties.\\ \tt L'altro giocatore ha inviato questo messaggio:\\ \tt \\ \tt \\ \tt "Sembra che stiamo avendo problemi a trovare una posizione libera per spostare il mio "A". Potremmo valutare di spostare altri oggetti per liberare spazio o trovare un'altra strategia per allineare gli oggetti identici. Che ne pensi? Dovremmo continuare a cercare una posizione per il mio "A" o provare un approccio diverso?"\\ \tt \\ \tt \\ \tt Qual è il tuo prossimo comando? 
	  } 
	   } 
	   } 
	 & \\ 
 

    \theutterance \stepcounter{utterance}  

    & \multicolumn{4}{p{0.6\linewidth}}{\cellcolor[rgb]{0.8,1,0.9}{%
	\makecell[{{p{\linewidth}}}]{% 
	  \tt {\tiny [A$\rangle$GM]}  
	 COMUNICA: Sì, sembra che stiamo avendo problemi a trovare una posizione libera per il tuo "A". Considerando che stiamo accumulando molte penalità, potrebbe essere meglio provare un approccio diverso per allineare gli oggetti identici. Una possibilità potrebbe essere di concentrarsi sugli oggetti che sono già vicini tra loro, come il tuo "S" e il mio "S", e cercare di allineare gli altri oggetti intorno a essi. Un'altra possibilità potrebbe essere di scambiare informazioni sulla disposizione degli oggetti sul nostro tavolo da gioco e cercare di trovare una configurazione che funzioni per entrambi. Che ne pensi di provare a scambiare informazioni sulla disposizione degli oggetti e vedere se possiamo trovare una soluzione che funzioni per entrambi? Inoltre, dovremmo anche valutare se è il momento di iniziare a pensare a come posizionare gli oggetti "T", "O", "D" e "M" per completare l'allineamento. 
	  } 
	   } 
	   } 
	 & & \\ 
 

    \theutterance \stepcounter{utterance}  

    & & & \multicolumn{2}{p{0.3\linewidth}}{\cellcolor[rgb]{0.95,0.95,0.95}{%
	\makecell[{{p{\linewidth}}}]{% 
	  \tt {\tiny [GM$|$GM]}  
	 COMUNICA: Sì, sembra che stiamo avendo problemi a trovare una posizione libera per il tuo "A". Considerando che stiamo accumulando molte penalità, potrebbe essere meglio provare un approccio diverso per allineare gli oggetti identici. Una possibilità potrebbe essere di concentrarsi sugli oggetti che sono già vicini tra loro, come il tuo "S" e il mio "S", e cercare di allineare gli altri oggetti intorno a essi. Un'altra possibilità potrebbe essere di scambiare informazioni sulla disposizione degli oggetti sul nostro tavolo da gioco e cercare di trovare una configurazione che funzioni per entrambi. Che ne pensi di provare a scambiare informazioni sulla disposizione degli oggetti e vedere se possiamo trovare una soluzione che funzioni per entrambi? Inoltre, dovremmo anche valutare se è il momento di iniziare a pensare a come posizionare gli oggetti "T", "O", "D" e "M" per completare l'allineamento. 
	  } 
	   } 
	   } 
	 & & \\ 
 

    \theutterance \stepcounter{utterance}  

    & & \multicolumn{4}{p{0.6\linewidth}}{\cellcolor[rgb]{0.9,0.9,0.9}{%
	\makecell[{{p{\linewidth}}}]{% 
	  \tt {\tiny [GM$\rangle$B]}  
	 Il tuo messaggio è stato inoltrato all'altro giocatore.\\ \tt \\ \tt \\ \tt You are currently playing round 9 of maximum 28.\\ \tt You have collectively accumulated 15 of 16 penalties.\\ \tt L'altro giocatore ha inviato questo messaggio:\\ \tt \\ \tt \\ \tt "Sì, sembra che stiamo avendo problemi a trovare una posizione libera per il tuo "A". Considerando che stiamo accumulando molte penalità, potrebbe essere meglio provare un approccio diverso per allineare gli oggetti identici. Una possibilità potrebbe essere di concentrarsi sugli oggetti che sono già vicini tra loro, come il tuo "S" e il mio "S", e cercare di allineare gli altri oggetti intorno a essi. Un'altra possibilità potrebbe essere di scambiare informazioni sulla disposizione degli oggetti sul nostro tavolo da gioco e cercare di trovare una configurazione che funzioni per entrambi. Che ne pensi di provare a scambiare informazioni sulla disposizione degli oggetti e vedere se possiamo trovare una soluzione che funzioni per entrambi? Inoltre, dovremmo anche valutare se è il momento di iniziare a pensare a come posizionare gli oggetti "T", "O", "D" e "M" per completare l'allineamento."\\ \tt \\ \tt \\ \tt Qual è il tuo prossimo comando? 
	  } 
	   } 
	   } 
	 & \\ 
 

    \theutterance \stepcounter{utterance}  

    & & & \multicolumn{4}{p{0.6\linewidth}}{\cellcolor[rgb]{1,0.85,0.72}{%
	\makecell[{{p{\linewidth}}}]{% 
	  \tt {\tiny [GM$\langle$B]}  
	 COMUNICA: Sono d'accordo che provare un approccio diverso potrebbe essere una buona idea. Scambiare informazioni sulla disposizione degli oggetti potrebbe aiutare a trovare una configurazione che funzioni per entrambi. Sono disposto a condividere la disposizione dei miei oggetti e ascoltare la tua. Inoltre, credo che valutare la posizione degli oggetti "T", "O", "D" e "M" sia importante per completare l'allineamento. Potresti iniziare a condividere la disposizione dei tuoi oggetti e io farò lo stesso? 
	  } 
	   } 
	   } 
	  \\ 
 

    \theutterance \stepcounter{utterance}  

    & & & \multicolumn{2}{p{0.3\linewidth}}{\cellcolor[rgb]{0.95,0.95,0.95}{%
	\makecell[{{p{\linewidth}}}]{% 
	  \tt {\tiny [GM$|$GM]}  
	 COMUNICA: Sono d'accordo che provare un approccio diverso potrebbe essere una buona idea. Scambiare informazioni sulla disposizione degli oggetti potrebbe aiutare a trovare una configurazione che funzioni per entrambi. Sono disposto a condividere la disposizione dei miei oggetti e ascoltare la tua. Inoltre, credo che valutare la posizione degli oggetti "T", "O", "D" e "M" sia importante per completare l'allineamento. Potresti iniziare a condividere la disposizione dei tuoi oggetti e io farò lo stesso? 
	  } 
	   } 
	   } 
	 & & \\ 
 

    \theutterance \stepcounter{utterance}  

    & & \multicolumn{4}{p{0.6\linewidth}}{\cellcolor[rgb]{0.9,0.9,0.9}{%
	\makecell[{{p{\linewidth}}}]{% 
	  \tt {\tiny [A$\langle$GM]}  
	 Il tuo messaggio è stato inoltrato all'altro giocatore.\\ \tt \\ \tt \\ \tt You are currently playing round 9 of maximum 28.\\ \tt You have collectively accumulated 15 of 16 penalties.\\ \tt L'altro giocatore ha inviato questo messaggio:\\ \tt \\ \tt \\ \tt "Sono d'accordo che provare un approccio diverso potrebbe essere una buona idea. Scambiare informazioni sulla disposizione degli oggetti potrebbe aiutare a trovare una configurazione che funzioni per entrambi. Sono disposto a condividere la disposizione dei miei oggetti e ascoltare la tua. Inoltre, credo che valutare la posizione degli oggetti "T", "O", "D" e "M" sia importante per completare l'allineamento. Potresti iniziare a condividere la disposizione dei tuoi oggetti e io farò lo stesso?"\\ \tt \\ \tt \\ \tt Qual è il tuo prossimo comando? 
	  } 
	   } 
	   } 
	 & \\ 
 

    \theutterance \stepcounter{utterance}  

    & \multicolumn{4}{p{0.6\linewidth}}{\cellcolor[rgb]{0.8,1,0.9}{%
	\makecell[{{p{\linewidth}}}]{% 
	  \tt {\tiny [A$\rangle$GM]}  
	 COMUNICA: Sì, sono d'accordo! Inizierò a condividere la disposizione dei miei oggetti. Il mio "P" è nella posizione (1,2), il mio "T" è nella posizione (7,7), il mio "O" è nella posizione (5,5), il mio "S" è nella posizione (3,3), il mio "D" è nella posizione (6,6), il mio "A" è nella posizione (3,3) accanto al mio "S", e il mio "M" è nella posizione (6,6) accanto al mio "D". Ora è il tuo turno! Per favore, condividi la disposizione dei tuoi oggetti in modo che possiamo lavorare insieme per trovare una configurazione che funzioni per entrambi. Inoltre, ricordiamoci che dobbiamo essere prudenti con le mosse per evitare di accumulare altre penalità. 
	  } 
	   } 
	   } 
	 & & \\ 
 

    \theutterance \stepcounter{utterance}  

    & & & \multicolumn{2}{p{0.3\linewidth}}{\cellcolor[rgb]{0.95,0.95,0.95}{%
	\makecell[{{p{\linewidth}}}]{% 
	  \tt {\tiny [GM$|$GM]}  
	 COMUNICA: Sì, sono d'accordo! Inizierò a condividere la disposizione dei miei oggetti. Il mio "P" è nella posizione (1,2), il mio "T" è nella posizione (7,7), il mio "O" è nella posizione (5,5), il mio "S" è nella posizione (3,3), il mio "D" è nella posizione (6,6), il mio "A" è nella posizione (3,3) accanto al mio "S", e il mio "M" è nella posizione (6,6) accanto al mio "D". Ora è il tuo turno! Per favore, condividi la disposizione dei tuoi oggetti in modo che possiamo lavorare insieme per trovare una configurazione che funzioni per entrambi. Inoltre, ricordiamoci che dobbiamo essere prudenti con le mosse per evitare di accumulare altre penalità. 
	  } 
	   } 
	   } 
	 & & \\ 
 

    \theutterance \stepcounter{utterance}  

    & & \multicolumn{4}{p{0.6\linewidth}}{\cellcolor[rgb]{0.9,0.9,0.9}{%
	\makecell[{{p{\linewidth}}}]{% 
	  \tt {\tiny [GM$\rangle$B]}  
	 Il tuo messaggio è stato inoltrato all'altro giocatore.\\ \tt \\ \tt \\ \tt You are currently playing round 10 of maximum 28.\\ \tt You have collectively accumulated 15 of 16 penalties.\\ \tt L'altro giocatore ha inviato questo messaggio:\\ \tt \\ \tt \\ \tt "Sì, sono d'accordo! Inizierò a condividere la disposizione dei miei oggetti. Il mio "P" è nella posizione (1,2), il mio "T" è nella posizione (7,7), il mio "O" è nella posizione (5,5), il mio "S" è nella posizione (3,3), il mio "D" è nella posizione (6,6), il mio "A" è nella posizione (3,3) accanto al mio "S", e il mio "M" è nella posizione (6,6) accanto al mio "D". Ora è il tuo turno! Per favore, condividi la disposizione dei tuoi oggetti in modo che possiamo lavorare insieme per trovare una configurazione che funzioni per entrambi. Inoltre, ricordiamoci che dobbiamo essere prudenti con le mosse per evitare di accumulare altre penalità."\\ \tt \\ \tt \\ \tt Qual è il tuo prossimo comando? 
	  } 
	   } 
	   } 
	 & \\ 
 

    \theutterance \stepcounter{utterance}  

    & & & \multicolumn{4}{p{0.6\linewidth}}{\cellcolor[rgb]{1,0.85,0.72}{%
	\makecell[{{p{\linewidth}}}]{% 
	  \tt {\tiny [GM$\langle$B]}  
	 COMUNICA: Grazie per aver condiviso la disposizione dei tuoi oggetti! Ecco la disposizione dei miei oggetti: il mio "P" è nella posizione (4,1), il mio "T" è nella posizione (2,2), il mio "O" è nella posizione (5,5), il mio "S" è nella posizione (4,5), il mio "D" è nella posizione (2,5), il mio "A" è nella posizione (1,7) e il mio "M" è nella posizione (3,3). Ora che abbiamo condiviso la disposizione dei nostri oggetti, possiamo lavorare insieme per trovare una configurazione che funzioni per entrambi. Potremmo iniziare a identificare gli oggetti che sono già vicini tra loro e trovare un modo per allinearli. Che ne pensi di iniziare con gli oggetti "O" e "M", che sembrano essere in posizioni simili sui nostri tavoli da gioco? 
	  } 
	   } 
	   } 
	  \\ 
 

    \theutterance \stepcounter{utterance}  

    & & & \multicolumn{2}{p{0.3\linewidth}}{\cellcolor[rgb]{0.95,0.95,0.95}{%
	\makecell[{{p{\linewidth}}}]{% 
	  \tt {\tiny [GM$|$GM]}  
	 COMUNICA: Grazie per aver condiviso la disposizione dei tuoi oggetti! Ecco la disposizione dei miei oggetti: il mio "P" è nella posizione (4,1), il mio "T" è nella posizione (2,2), il mio "O" è nella posizione (5,5), il mio "S" è nella posizione (4,5), il mio "D" è nella posizione (2,5), il mio "A" è nella posizione (1,7) e il mio "M" è nella posizione (3,3). Ora che abbiamo condiviso la disposizione dei nostri oggetti, possiamo lavorare insieme per trovare una configurazione che funzioni per entrambi. Potremmo iniziare a identificare gli oggetti che sono già vicini tra loro e trovare un modo per allinearli. Che ne pensi di iniziare con gli oggetti "O" e "M", che sembrano essere in posizioni simili sui nostri tavoli da gioco? 
	  } 
	   } 
	   } 
	 & & \\ 
 

    \theutterance \stepcounter{utterance}  

    & & \multicolumn{4}{p{0.6\linewidth}}{\cellcolor[rgb]{0.9,0.9,0.9}{%
	\makecell[{{p{\linewidth}}}]{% 
	  \tt {\tiny [A$\langle$GM]}  
	 Il tuo messaggio è stato inoltrato all'altro giocatore.\\ \tt \\ \tt \\ \tt You are currently playing round 10 of maximum 28.\\ \tt You have collectively accumulated 15 of 16 penalties.\\ \tt L'altro giocatore ha inviato questo messaggio:\\ \tt \\ \tt \\ \tt "Grazie per aver condiviso la disposizione dei tuoi oggetti! Ecco la disposizione dei miei oggetti: il mio "P" è nella posizione (4,1), il mio "T" è nella posizione (2,2), il mio "O" è nella posizione (5,5), il mio "S" è nella posizione (4,5), il mio "D" è nella posizione (2,5), il mio "A" è nella posizione (1,7) e il mio "M" è nella posizione (3,3). Ora che abbiamo condiviso la disposizione dei nostri oggetti, possiamo lavorare insieme per trovare una configurazione che funzioni per entrambi. Potremmo iniziare a identificare gli oggetti che sono già vicini tra loro e trovare un modo per allinearli. Che ne pensi di iniziare con gli oggetti "O" e "M", che sembrano essere in posizioni simili sui nostri tavoli da gioco?"\\ \tt \\ \tt \\ \tt Qual è il tuo prossimo comando? 
	  } 
	   } 
	   } 
	 & \\ 
 

    \theutterance \stepcounter{utterance}  

    & \multicolumn{4}{p{0.6\linewidth}}{\cellcolor[rgb]{0.8,1,0.9}{%
	\makecell[{{p{\linewidth}}}]{% 
	  \tt {\tiny [A$\rangle$GM]}  
	 COMUNICA: Sì, iniziare con gli oggetti "O" e "M" sembra una buona idea, considerando che sono già in posizioni simili sui nostri tavoli da gioco. Notiamo che il mio "O" è nella posizione (5,5) e il tuo "O" è nella stessa posizione (5,5), quindi sono già allineati! Questo è un buon punto di partenza. Per quanto riguarda il mio "M" e il tuo "M", il mio è nella posizione (6,6) e il tuo è nella posizione (3,3), quindi non sono allineati. Potremmo provare a spostare il mio "M" o il tuo "M" per allinearli. Tuttavia, dobbiamo essere cauti e valutare le conseguenze di ogni mossa per evitare di accumulare altre penalità. Che ne pensi di analizzare le possibilità di spostamento per il mio "M" e il tuo "M" e discuterne le opzioni prima di prendere una decisione? 
	  } 
	   } 
	   } 
	 & & \\ 
 

    \theutterance \stepcounter{utterance}  

    & & & \multicolumn{2}{p{0.3\linewidth}}{\cellcolor[rgb]{0.95,0.95,0.95}{%
	\makecell[{{p{\linewidth}}}]{% 
	  \tt {\tiny [GM$|$GM]}  
	 COMUNICA: Sì, iniziare con gli oggetti "O" e "M" sembra una buona idea, considerando che sono già in posizioni simili sui nostri tavoli da gioco. Notiamo che il mio "O" è nella posizione (5,5) e il tuo "O" è nella stessa posizione (5,5), quindi sono già allineati! Questo è un buon punto di partenza. Per quanto riguarda il mio "M" e il tuo "M", il mio è nella posizione (6,6) e il tuo è nella posizione (3,3), quindi non sono allineati. Potremmo provare a spostare il mio "M" o il tuo "M" per allinearli. Tuttavia, dobbiamo essere cauti e valutare le conseguenze di ogni mossa per evitare di accumulare altre penalità. Che ne pensi di analizzare le possibilità di spostamento per il mio "M" e il tuo "M" e discuterne le opzioni prima di prendere una decisione? 
	  } 
	   } 
	   } 
	 & & \\ 
 

    \theutterance \stepcounter{utterance}  

    & & \multicolumn{4}{p{0.6\linewidth}}{\cellcolor[rgb]{0.9,0.9,0.9}{%
	\makecell[{{p{\linewidth}}}]{% 
	  \tt {\tiny [GM$\rangle$B]}  
	 Il tuo messaggio è stato inoltrato all'altro giocatore.\\ \tt \\ \tt \\ \tt You are currently playing round 11 of maximum 28.\\ \tt You have collectively accumulated 15 of 16 penalties.\\ \tt L'altro giocatore ha inviato questo messaggio:\\ \tt \\ \tt \\ \tt "Sì, iniziare con gli oggetti "O" e "M" sembra una buona idea, considerando che sono già in posizioni simili sui nostri tavoli da gioco. Notiamo che il mio "O" è nella posizione (5,5) e il tuo "O" è nella stessa posizione (5,5), quindi sono già allineati! Questo è un buon punto di partenza. Per quanto riguarda il mio "M" e il tuo "M", il mio è nella posizione (6,6) e il tuo è nella posizione (3,3), quindi non sono allineati. Potremmo provare a spostare il mio "M" o il tuo "M" per allinearli. Tuttavia, dobbiamo essere cauti e valutare le conseguenze di ogni mossa per evitare di accumulare altre penalità. Che ne pensi di analizzare le possibilità di spostamento per il mio "M" e il tuo "M" e discuterne le opzioni prima di prendere una decisione?"\\ \tt \\ \tt \\ \tt Qual è il tuo prossimo comando? 
	  } 
	   } 
	   } 
	 & \\ 
 

    \theutterance \stepcounter{utterance}  

    & & & \multicolumn{4}{p{0.6\linewidth}}{\cellcolor[rgb]{1,0.85,0.72}{%
	\makecell[{{p{\linewidth}}}]{% 
	  \tt {\tiny [GM$\langle$B]}  
	 COMUNICA: Sono d'accordo che analizzare le possibilità di spostamento per il tuo "M" e il mio "M" sia una buona idea. Considerando che il mio "M" è nella posizione (3,3) e il tuo "M" è nella posizione (6,6), potremmo valutare le opzioni per spostare uno di essi nella stessa posizione dell'altro. Tuttavia, dobbiamo anche considerare la disposizione degli altri oggetti intorno a loro per evitare di creare conflitti o accumulare altre penalità. Potresti Farmi una lista delle posizioni adiacenti al tuo "M" e al mio "M" per capire quali sono le opzioni disponibili per lo spostamento? 
	  } 
	   } 
	   } 
	  \\ 
 

    \theutterance \stepcounter{utterance}  

    & & & \multicolumn{2}{p{0.3\linewidth}}{\cellcolor[rgb]{0.95,0.95,0.95}{%
	\makecell[{{p{\linewidth}}}]{% 
	  \tt {\tiny [GM$|$GM]}  
	 COMUNICA: Sono d'accordo che analizzare le possibilità di spostamento per il tuo "M" e il mio "M" sia una buona idea. Considerando che il mio "M" è nella posizione (3,3) e il tuo "M" è nella posizione (6,6), potremmo valutare le opzioni per spostare uno di essi nella stessa posizione dell'altro. Tuttavia, dobbiamo anche considerare la disposizione degli altri oggetti intorno a loro per evitare di creare conflitti o accumulare altre penalità. Potresti Farmi una lista delle posizioni adiacenti al tuo "M" e al mio "M" per capire quali sono le opzioni disponibili per lo spostamento? 
	  } 
	   } 
	   } 
	 & & \\ 
 

    \theutterance \stepcounter{utterance}  

    & & \multicolumn{4}{p{0.6\linewidth}}{\cellcolor[rgb]{0.9,0.9,0.9}{%
	\makecell[{{p{\linewidth}}}]{% 
	  \tt {\tiny [A$\langle$GM]}  
	 Il tuo messaggio è stato inoltrato all'altro giocatore.\\ \tt \\ \tt \\ \tt You are currently playing round 11 of maximum 28.\\ \tt You have collectively accumulated 15 of 16 penalties.\\ \tt L'altro giocatore ha inviato questo messaggio:\\ \tt \\ \tt \\ \tt "Sono d'accordo che analizzare le possibilità di spostamento per il tuo "M" e il mio "M" sia una buona idea. Considerando che il mio "M" è nella posizione (3,3) e il tuo "M" è nella posizione (6,6), potremmo valutare le opzioni per spostare uno di essi nella stessa posizione dell'altro. Tuttavia, dobbiamo anche considerare la disposizione degli altri oggetti intorno a loro per evitare di creare conflitti o accumulare altre penalità. Potresti Farmi una lista delle posizioni adiacenti al tuo "M" e al mio "M" per capire quali sono le opzioni disponibili per lo spostamento?"\\ \tt \\ \tt \\ \tt Qual è il tuo prossimo comando? 
	  } 
	   } 
	   } 
	 & \\ 
 

    \theutterance \stepcounter{utterance}  

    & \multicolumn{4}{p{0.6\linewidth}}{\cellcolor[rgb]{0.8,1,0.9}{%
	\makecell[{{p{\linewidth}}}]{% 
	  \tt {\tiny [A$\rangle$GM]}  
	 COMUNICA: Sì, certo! Le posizioni adiacenti al mio "M" (6,6) sono: (5,6), (6,5), (6,7), (7,6), (5,5) e (7,5). Le posizioni adiacenti al tuo "M" (3,3) sono: (2,3), (3,2), (3,4), (4,3), (2,2) e (4,2). Tuttavia, dobbiamo anche considerare che alcune di queste posizioni potrebbero essere occupate da altri oggetti, quindi è importante valutare la disposizione degli oggetti intorno ai nostri "M" prima di prendere una decisione. Inoltre, potremmo anche considerare di spostare altri oggetti per creare spazio o allineare gli oggetti identici. Che ne pensi di valutare la disposizione degli oggetti intorno ai nostri "M" e discutere le opzioni di spostamento in base a questa informazione? 
	  } 
	   } 
	   } 
	 & & \\ 
 

    \theutterance \stepcounter{utterance}  

    & & & \multicolumn{2}{p{0.3\linewidth}}{\cellcolor[rgb]{0.95,0.95,0.95}{%
	\makecell[{{p{\linewidth}}}]{% 
	  \tt {\tiny [GM$|$GM]}  
	 COMUNICA: Sì, certo! Le posizioni adiacenti al mio "M" (6,6) sono: (5,6), (6,5), (6,7), (7,6), (5,5) e (7,5). Le posizioni adiacenti al tuo "M" (3,3) sono: (2,3), (3,2), (3,4), (4,3), (2,2) e (4,2). Tuttavia, dobbiamo anche considerare che alcune di queste posizioni potrebbero essere occupate da altri oggetti, quindi è importante valutare la disposizione degli oggetti intorno ai nostri "M" prima di prendere una decisione. Inoltre, potremmo anche considerare di spostare altri oggetti per creare spazio o allineare gli oggetti identici. Che ne pensi di valutare la disposizione degli oggetti intorno ai nostri "M" e discutere le opzioni di spostamento in base a questa informazione? 
	  } 
	   } 
	   } 
	 & & \\ 
 

    \theutterance \stepcounter{utterance}  

    & & \multicolumn{4}{p{0.6\linewidth}}{\cellcolor[rgb]{0.9,0.9,0.9}{%
	\makecell[{{p{\linewidth}}}]{% 
	  \tt {\tiny [GM$\rangle$B]}  
	 Il tuo messaggio è stato inoltrato all'altro giocatore.\\ \tt \\ \tt \\ \tt You are currently playing round 12 of maximum 28.\\ \tt You have collectively accumulated 15 of 16 penalties.\\ \tt L'altro giocatore ha inviato questo messaggio:\\ \tt \\ \tt \\ \tt "Sì, certo! Le posizioni adiacenti al mio "M" (6,6) sono: (5,6), (6,5), (6,7), (7,6), (5,5) e (7,5). Le posizioni adiacenti al tuo "M" (3,3) sono: (2,3), (3,2), (3,4), (4,3), (2,2) e (4,2). Tuttavia, dobbiamo anche considerare che alcune di queste posizioni potrebbero essere occupate da altri oggetti, quindi è importante valutare la disposizione degli oggetti intorno ai nostri "M" prima di prendere una decisione. Inoltre, potremmo anche considerare di spostare altri oggetti per creare spazio o allineare gli oggetti identici. Che ne pensi di valutare la disposizione degli oggetti intorno ai nostri "M" e discutere le opzioni di spostamento in base a questa informazione?"\\ \tt \\ \tt \\ \tt Qual è il tuo prossimo comando? 
	  } 
	   } 
	   } 
	 & \\ 
 

    \theutterance \stepcounter{utterance}  

    & & & \multicolumn{4}{p{0.6\linewidth}}{\cellcolor[rgb]{1,0.85,0.72}{%
	\makecell[{{p{\linewidth}}}]{% 
	  \tt {\tiny [GM$\langle$B]}  
	 COMUNICA: Sono d'accordo che valutare la disposizione degli oggetti intorno ai nostri "M" sia importante. Considerando le posizioni adiacenti al mio "M" e al tuo "M", potremmo iniziare a identificare le posizioni che sono libere e quelle che sono occupate da altri oggetti. Potresti condividere la disposizione degli oggetti intorno al tuo "M" e io farò lo stesso per il mio "M"? In questo modo, potremmo avere una visione più chiara delle opzioni di spostamento disponibili e prendere una decisione informata. Inoltre, vorrei anche discutere la possibilità di spostare altri oggetti per creare spazio o allineare gli oggetti identici, come hai suggerito. 
	  } 
	   } 
	   } 
	  \\ 
 

    \theutterance \stepcounter{utterance}  

    & & & \multicolumn{2}{p{0.3\linewidth}}{\cellcolor[rgb]{0.95,0.95,0.95}{%
	\makecell[{{p{\linewidth}}}]{% 
	  \tt {\tiny [GM$|$GM]}  
	 COMUNICA: Sono d'accordo che valutare la disposizione degli oggetti intorno ai nostri "M" sia importante. Considerando le posizioni adiacenti al mio "M" e al tuo "M", potremmo iniziare a identificare le posizioni che sono libere e quelle che sono occupate da altri oggetti. Potresti condividere la disposizione degli oggetti intorno al tuo "M" e io farò lo stesso per il mio "M"? In questo modo, potremmo avere una visione più chiara delle opzioni di spostamento disponibili e prendere una decisione informata. Inoltre, vorrei anche discutere la possibilità di spostare altri oggetti per creare spazio o allineare gli oggetti identici, come hai suggerito. 
	  } 
	   } 
	   } 
	 & & \\ 
 

    \theutterance \stepcounter{utterance}  

    & & \multicolumn{4}{p{0.6\linewidth}}{\cellcolor[rgb]{0.9,0.9,0.9}{%
	\makecell[{{p{\linewidth}}}]{% 
	  \tt {\tiny [A$\langle$GM]}  
	 Il tuo messaggio è stato inoltrato all'altro giocatore.\\ \tt \\ \tt \\ \tt You are currently playing round 12 of maximum 28.\\ \tt You have collectively accumulated 15 of 16 penalties.\\ \tt L'altro giocatore ha inviato questo messaggio:\\ \tt \\ \tt \\ \tt "Sono d'accordo che valutare la disposizione degli oggetti intorno ai nostri "M" sia importante. Considerando le posizioni adiacenti al mio "M" e al tuo "M", potremmo iniziare a identificare le posizioni che sono libere e quelle che sono occupate da altri oggetti. Potresti condividere la disposizione degli oggetti intorno al tuo "M" e io farò lo stesso per il mio "M"? In questo modo, potremmo avere una visione più chiara delle opzioni di spostamento disponibili e prendere una decisione informata. Inoltre, vorrei anche discutere la possibilità di spostare altri oggetti per creare spazio o allineare gli oggetti identici, come hai suggerito."\\ \tt \\ \tt \\ \tt Qual è il tuo prossimo comando? 
	  } 
	   } 
	   } 
	 & \\ 
 

    \theutterance \stepcounter{utterance}  

    & \multicolumn{4}{p{0.6\linewidth}}{\cellcolor[rgb]{0.8,1,0.9}{%
	\makecell[{{p{\linewidth}}}]{% 
	  \tt {\tiny [A$\rangle$GM]}  
	 COMUNICA: Sì, certo! Intorno al mio "M" (6,6), ci sono i seguenti oggetti: il mio "D" è nella posizione (6,6) accanto al mio "M", quindi occupa la stessa cella. Le celle adiacenti (5,6), (6,5), (6,7) e (7,6) sono libere, ma la cella (5,5) è occupata dal mio "O". La cella (7,5) è libera. Per quanto riguarda gli altri oggetti, il mio "P" è nella posizione (1,2), il mio "T" è nella posizione (7,7), il mio "S" è nella posizione (3,3) e il mio "A" è nella posizione (3,3) accanto al mio "S". Ora è il tuo turno! Per favore, condividi la disposizione degli oggetti intorno al tuo "M" (3,3) in modo che possiamo valutare le opzioni di spostamento disponibili e prendere una decisione informata. Inoltre, sono ansioso di discutere la possibilità di spostare altri oggetti per creare spazio o allineare gli oggetti identici. 
	  } 
	   } 
	   } 
	 & & \\ 
 

    \theutterance \stepcounter{utterance}  

    & & & \multicolumn{2}{p{0.3\linewidth}}{\cellcolor[rgb]{0.95,0.95,0.95}{%
	\makecell[{{p{\linewidth}}}]{% 
	  \tt {\tiny [GM$|$GM]}  
	 COMUNICA: Sì, certo! Intorno al mio "M" (6,6), ci sono i seguenti oggetti: il mio "D" è nella posizione (6,6) accanto al mio "M", quindi occupa la stessa cella. Le celle adiacenti (5,6), (6,5), (6,7) e (7,6) sono libere, ma la cella (5,5) è occupata dal mio "O". La cella (7,5) è libera. Per quanto riguarda gli altri oggetti, il mio "P" è nella posizione (1,2), il mio "T" è nella posizione (7,7), il mio "S" è nella posizione (3,3) e il mio "A" è nella posizione (3,3) accanto al mio "S". Ora è il tuo turno! Per favore, condividi la disposizione degli oggetti intorno al tuo "M" (3,3) in modo che possiamo valutare le opzioni di spostamento disponibili e prendere una decisione informata. Inoltre, sono ansioso di discutere la possibilità di spostare altri oggetti per creare spazio o allineare gli oggetti identici. 
	  } 
	   } 
	   } 
	 & & \\ 
 

    \theutterance \stepcounter{utterance}  

    & & \multicolumn{4}{p{0.6\linewidth}}{\cellcolor[rgb]{0.9,0.9,0.9}{%
	\makecell[{{p{\linewidth}}}]{% 
	  \tt {\tiny [GM$\rangle$B]}  
	 Il tuo messaggio è stato inoltrato all'altro giocatore.\\ \tt \\ \tt \\ \tt You are currently playing round 13 of maximum 28.\\ \tt You have collectively accumulated 15 of 16 penalties.\\ \tt L'altro giocatore ha inviato questo messaggio:\\ \tt \\ \tt \\ \tt "Sì, certo! Intorno al mio "M" (6,6), ci sono i seguenti oggetti: il mio "D" è nella posizione (6,6) accanto al mio "M", quindi occupa la stessa cella. Le celle adiacenti (5,6), (6,5), (6,7) e (7,6) sono libere, ma la cella (5,5) è occupata dal mio "O". La cella (7,5) è libera. Per quanto riguarda gli altri oggetti, il mio "P" è nella posizione (1,2), il mio "T" è nella posizione (7,7), il mio "S" è nella posizione (3,3) e il mio "A" è nella posizione (3,3) accanto al mio "S". Ora è il tuo turno! Per favore, condividi la disposizione degli oggetti intorno al tuo "M" (3,3) in modo che possiamo valutare le opzioni di spostamento disponibili e prendere una decisione informata. Inoltre, sono ansioso di discutere la possibilità di spostare altri oggetti per creare spazio o allineare gli oggetti identici."\\ \tt \\ \tt \\ \tt Qual è il tuo prossimo comando? 
	  } 
	   } 
	   } 
	 & \\ 
 

    \theutterance \stepcounter{utterance}  

    & & & \multicolumn{4}{p{0.6\linewidth}}{\cellcolor[rgb]{1,0.85,0.72}{%
	\makecell[{{p{\linewidth}}}]{% 
	  \tt {\tiny [GM$\langle$B]}  
	 COMUNICA: Grazie per la descrizione degli oggetti intorno al tuo "M"! Intorno al mio "M" (3,3), ci sono i seguenti oggetti: il mio "P" è nella posizione (4,1), il mio "T" è nella posizione (2,2), il mio "S" è nella posizione (4,5), il mio "D" è nella posizione (2,5) e il mio "A" è nella posizione (1,7). Le celle adiacenti al mio "M" (3,3) sono: (2,3), (3,2), (3,4), (4,3), (2,2) e (4,2). La cella (2,3) è libera, la cella (3,2) è libera, la cella (3,4) è occupata da una linea orizzontale, la cella (4,3) è libera, la cella (2,2) è occupata dal mio "T" e la cella (4,2) è libera. Ora che abbiamo una visione completa degli oggetti intorno ai nostri "M", possiamo valutare le opzioni di spostamento disponibili e prendere una decisione informata. Sembra che ci siano alcune celle libere intorno ai nostri "M" che potremmo utilizzare per spostare gli oggetti e allinearli. Che ne pensi di iniziare a valutare le opzioni di spostamento per il mio "M" e il tuo "M"? 
	  } 
	   } 
	   } 
	  \\ 
 

    \theutterance \stepcounter{utterance}  

    & & & \multicolumn{2}{p{0.3\linewidth}}{\cellcolor[rgb]{0.95,0.95,0.95}{%
	\makecell[{{p{\linewidth}}}]{% 
	  \tt {\tiny [GM$|$GM]}  
	 COMUNICA: Grazie per la descrizione degli oggetti intorno al tuo "M"! Intorno al mio "M" (3,3), ci sono i seguenti oggetti: il mio "P" è nella posizione (4,1), il mio "T" è nella posizione (2,2), il mio "S" è nella posizione (4,5), il mio "D" è nella posizione (2,5) e il mio "A" è nella posizione (1,7). Le celle adiacenti al mio "M" (3,3) sono: (2,3), (3,2), (3,4), (4,3), (2,2) e (4,2). La cella (2,3) è libera, la cella (3,2) è libera, la cella (3,4) è occupata da una linea orizzontale, la cella (4,3) è libera, la cella (2,2) è occupata dal mio "T" e la cella (4,2) è libera. Ora che abbiamo una visione completa degli oggetti intorno ai nostri "M", possiamo valutare le opzioni di spostamento disponibili e prendere una decisione informata. Sembra che ci siano alcune celle libere intorno ai nostri "M" che potremmo utilizzare per spostare gli oggetti e allinearli. Che ne pensi di iniziare a valutare le opzioni di spostamento per il mio "M" e il tuo "M"? 
	  } 
	   } 
	   } 
	 & & \\ 
 

    \theutterance \stepcounter{utterance}  

    & & \multicolumn{4}{p{0.6\linewidth}}{\cellcolor[rgb]{0.9,0.9,0.9}{%
	\makecell[{{p{\linewidth}}}]{% 
	  \tt {\tiny [A$\langle$GM]}  
	 Il tuo messaggio è stato inoltrato all'altro giocatore.\\ \tt \\ \tt \\ \tt You are currently playing round 13 of maximum 28.\\ \tt You have collectively accumulated 15 of 16 penalties.\\ \tt L'altro giocatore ha inviato questo messaggio:\\ \tt \\ \tt \\ \tt "Grazie per la descrizione degli oggetti intorno al tuo "M"! Intorno al mio "M" (3,3), ci sono i seguenti oggetti: il mio "P" è nella posizione (4,1), il mio "T" è nella posizione (2,2), il mio "S" è nella posizione (4,5), il mio "D" è nella posizione (2,5) e il mio "A" è nella posizione (1,7). Le celle adiacenti al mio "M" (3,3) sono: (2,3), (3,2), (3,4), (4,3), (2,2) e (4,2). La cella (2,3) è libera, la cella (3,2) è libera, la cella (3,4) è occupata da una linea orizzontale, la cella (4,3) è libera, la cella (2,2) è occupata dal mio "T" e la cella (4,2) è libera. Ora che abbiamo una visione completa degli oggetti intorno ai nostri "M", possiamo valutare le opzioni di spostamento disponibili e prendere una decisione informata. Sembra che ci siano alcune celle libere intorno ai nostri "M" che potremmo utilizzare per spostare gli oggetti e allinearli. Che ne pensi di iniziare a valutare le opzioni di spostamento per il mio "M" e il tuo "M"?"\\ \tt \\ \tt \\ \tt Qual è il tuo prossimo comando? 
	  } 
	   } 
	   } 
	 & \\ 
 

    \theutterance \stepcounter{utterance}  

    & \multicolumn{4}{p{0.6\linewidth}}{\cellcolor[rgb]{0.8,1,0.9}{%
	\makecell[{{p{\linewidth}}}]{% 
	  \tt {\tiny [A$\rangle$GM]}  
	 COMUNICA: Sì, sono d'accordo! Ora che abbiamo una visione completa degli oggetti intorno ai nostri "M", possiamo valutare le opzioni di spostamento disponibili. Una possibilità potrebbe essere di spostare il mio "M" nella posizione (4,3), che è libera, per allineararlo con il tuo "M". Tuttavia, dobbiamo considerare gli oggetti che sono già presenti nella posizione (4,3) o vicino a essa. Un'altra possibilità potrebbe essere di spostare il tuo "M" nella posizione (6,6), che è occupata dal mio "M" e dal mio "D", quindi dovremmo spostare uno di questi oggetti per fare spazio. Un'altra opzione potrebbe essere di spostare il tuo "M" nella posizione (5,5), che è occupata dal mio "O", quindi dovremmo spostare il mio "O" per fare spazio. Che ne pensi di valutare queste opzioni e discuterne le possibilità e le conseguenze? Inoltre, dovremmo anche considerare di spostare altri oggetti per creare spazio o allineare gli oggetti identici, come ad esempio spostare il mio "D" o il tuo "D" per fare spazio per il tuo "M" o il mio "M". 
	  } 
	   } 
	   } 
	 & & \\ 
 

    \theutterance \stepcounter{utterance}  

    & & & \multicolumn{2}{p{0.3\linewidth}}{\cellcolor[rgb]{0.95,0.95,0.95}{%
	\makecell[{{p{\linewidth}}}]{% 
	  \tt {\tiny [GM$|$GM]}  
	 COMUNICA: Sì, sono d'accordo! Ora che abbiamo una visione completa degli oggetti intorno ai nostri "M", possiamo valutare le opzioni di spostamento disponibili. Una possibilità potrebbe essere di spostare il mio "M" nella posizione (4,3), che è libera, per allineararlo con il tuo "M". Tuttavia, dobbiamo considerare gli oggetti che sono già presenti nella posizione (4,3) o vicino a essa. Un'altra possibilità potrebbe essere di spostare il tuo "M" nella posizione (6,6), che è occupata dal mio "M" e dal mio "D", quindi dovremmo spostare uno di questi oggetti per fare spazio. Un'altra opzione potrebbe essere di spostare il tuo "M" nella posizione (5,5), che è occupata dal mio "O", quindi dovremmo spostare il mio "O" per fare spazio. Che ne pensi di valutare queste opzioni e discuterne le possibilità e le conseguenze? Inoltre, dovremmo anche considerare di spostare altri oggetti per creare spazio o allineare gli oggetti identici, come ad esempio spostare il mio "D" o il tuo "D" per fare spazio per il tuo "M" o il mio "M". 
	  } 
	   } 
	   } 
	 & & \\ 
 

    \theutterance \stepcounter{utterance}  

    & & \multicolumn{4}{p{0.6\linewidth}}{\cellcolor[rgb]{0.9,0.9,0.9}{%
	\makecell[{{p{\linewidth}}}]{% 
	  \tt {\tiny [GM$\rangle$B]}  
	 Il tuo messaggio è stato inoltrato all'altro giocatore.\\ \tt \\ \tt \\ \tt You are currently playing round 14 of maximum 28.\\ \tt You have collectively accumulated 15 of 16 penalties.\\ \tt L'altro giocatore ha inviato questo messaggio:\\ \tt \\ \tt \\ \tt "Sì, sono d'accordo! Ora che abbiamo una visione completa degli oggetti intorno ai nostri "M", possiamo valutare le opzioni di spostamento disponibili. Una possibilità potrebbe essere di spostare il mio "M" nella posizione (4,3), che è libera, per allineararlo con il tuo "M". Tuttavia, dobbiamo considerare gli oggetti che sono già presenti nella posizione (4,3) o vicino a essa. Un'altra possibilità potrebbe essere di spostare il tuo "M" nella posizione (6,6), che è occupata dal mio "M" e dal mio "D", quindi dovremmo spostare uno di questi oggetti per fare spazio. Un'altra opzione potrebbe essere di spostare il tuo "M" nella posizione (5,5), che è occupata dal mio "O", quindi dovremmo spostare il mio "O" per fare spazio. Che ne pensi di valutare queste opzioni e discuterne le possibilità e le conseguenze? Inoltre, dovremmo anche considerare di spostare altri oggetti per creare spazio o allineare gli oggetti identici, come ad esempio spostare il mio "D" o il tuo "D" per fare spazio per il tuo "M" o il mio "M"."\\ \tt \\ \tt \\ \tt Qual è il tuo prossimo comando? 
	  } 
	   } 
	   } 
	 & \\ 
 

    \theutterance \stepcounter{utterance}  

    & & & \multicolumn{4}{p{0.6\linewidth}}{\cellcolor[rgb]{1,0.85,0.72}{%
	\makecell[{{p{\linewidth}}}]{% 
	  \tt {\tiny [GM$\langle$B]}  
	 COMUNICA: Sono d'accordo che valutare le opzioni di spostamento disponibili sia importante. Considerando le possibilità che hai menzionato, credo che spostare il mio "M" nella posizione (5,5) potrebbe essere una buona opzione, ma come hai detto, dovremmo spostare il tuo "O" per fare spazio. Un'altra opzione potrebbe essere di spostare il tuo "D" per fare spazio per il mio "M" nella posizione (6,6). Tuttavia, dobbiamo anche considerare le conseguenze di ogni mossa e assicurarci di non creare altri problemi. Potremmo iniziare a valutare le conseguenze di ogni opzione e discuterne le possibilità? Inoltre, credo che anche spostare il mio "D" potrebbe essere una buona opzione per creare spazio per il tuo "M" o il mio "M". Che ne pensi di valutare anche questa possibilità? 
	  } 
	   } 
	   } 
	  \\ 
 

    \theutterance \stepcounter{utterance}  

    & & & \multicolumn{2}{p{0.3\linewidth}}{\cellcolor[rgb]{0.95,0.95,0.95}{%
	\makecell[{{p{\linewidth}}}]{% 
	  \tt {\tiny [GM$|$GM]}  
	 COMUNICA: Sono d'accordo che valutare le opzioni di spostamento disponibili sia importante. Considerando le possibilità che hai menzionato, credo che spostare il mio "M" nella posizione (5,5) potrebbe essere una buona opzione, ma come hai detto, dovremmo spostare il tuo "O" per fare spazio. Un'altra opzione potrebbe essere di spostare il tuo "D" per fare spazio per il mio "M" nella posizione (6,6). Tuttavia, dobbiamo anche considerare le conseguenze di ogni mossa e assicurarci di non creare altri problemi. Potremmo iniziare a valutare le conseguenze di ogni opzione e discuterne le possibilità? Inoltre, credo che anche spostare il mio "D" potrebbe essere una buona opzione per creare spazio per il tuo "M" o il mio "M". Che ne pensi di valutare anche questa possibilità? 
	  } 
	   } 
	   } 
	 & & \\ 
 

    \theutterance \stepcounter{utterance}  

    & & \multicolumn{4}{p{0.6\linewidth}}{\cellcolor[rgb]{0.9,0.9,0.9}{%
	\makecell[{{p{\linewidth}}}]{% 
	  \tt {\tiny [A$\langle$GM]}  
	 Il tuo messaggio è stato inoltrato all'altro giocatore.\\ \tt \\ \tt \\ \tt You are currently playing round 14 of maximum 28.\\ \tt You have collectively accumulated 15 of 16 penalties.\\ \tt L'altro giocatore ha inviato questo messaggio:\\ \tt \\ \tt \\ \tt "Sono d'accordo che valutare le opzioni di spostamento disponibili sia importante. Considerando le possibilità che hai menzionato, credo che spostare il mio "M" nella posizione (5,5) potrebbe essere una buona opzione, ma come hai detto, dovremmo spostare il tuo "O" per fare spazio. Un'altra opzione potrebbe essere di spostare il tuo "D" per fare spazio per il mio "M" nella posizione (6,6). Tuttavia, dobbiamo anche considerare le conseguenze di ogni mossa e assicurarci di non creare altri problemi. Potremmo iniziare a valutare le conseguenze di ogni opzione e discuterne le possibilità? Inoltre, credo che anche spostare il mio "D" potrebbe essere una buona opzione per creare spazio per il tuo "M" o il mio "M". Che ne pensi di valutare anche questa possibilità?"\\ \tt \\ \tt \\ \tt Qual è il tuo prossimo comando? 
	  } 
	   } 
	   } 
	 & \\ 
 

    \theutterance \stepcounter{utterance}  

    & \multicolumn{4}{p{0.6\linewidth}}{\cellcolor[rgb]{0.8,1,0.9}{%
	\makecell[{{p{\linewidth}}}]{% 
	  \tt {\tiny [A$\rangle$GM]}  
	 COMUNICA: Sì, sono d'accordo! Valutare le conseguenze di ogni opzione è fondamentale per evitare di creare altri problemi. Spostare il mio "O" dalla posizione (5,5) per fare spazio per il tuo "M" potrebbe essere una buona opzione, ma dobbiamo considerare dove spostare il mio "O" e se ciò creerà altri problemi. Spostare il mio "D" dalla posizione (6,6) per fare spazio per il tuo "M" potrebbe essere un'altra opzione, ma dobbiamo considerare se ciò creerà spazio sufficiente per il tuo "M" e se ci saranno altri oggetti che potrebbero ostacolare lo spostamento. Spostare il tuo "D" dalla posizione (2,5) potrebbe essere un'altra opzione per creare spazio per il mio "M" o il tuo "M", ma dobbiamo considerare dove spostare il tuo "D" e se ciò creerà altri problemi. Che ne pensi di creare una lista delle opzioni di spostamento possibili e delle loro conseguenze per ogni oggetto, in modo da poter valutare meglio le possibilità e prendere una decisione informata? Inoltre, dovremmo anche considerare il numero di penalità che abbiamo accumulato e cercare di minimizzare il rischio di accumularne altre. 
	  } 
	   } 
	   } 
	 & & \\ 
 

    \theutterance \stepcounter{utterance}  

    & & & \multicolumn{2}{p{0.3\linewidth}}{\cellcolor[rgb]{0.95,0.95,0.95}{%
	\makecell[{{p{\linewidth}}}]{% 
	  \tt {\tiny [GM$|$GM]}  
	 COMUNICA: Sì, sono d'accordo! Valutare le conseguenze di ogni opzione è fondamentale per evitare di creare altri problemi. Spostare il mio "O" dalla posizione (5,5) per fare spazio per il tuo "M" potrebbe essere una buona opzione, ma dobbiamo considerare dove spostare il mio "O" e se ciò creerà altri problemi. Spostare il mio "D" dalla posizione (6,6) per fare spazio per il tuo "M" potrebbe essere un'altra opzione, ma dobbiamo considerare se ciò creerà spazio sufficiente per il tuo "M" e se ci saranno altri oggetti che potrebbero ostacolare lo spostamento. Spostare il tuo "D" dalla posizione (2,5) potrebbe essere un'altra opzione per creare spazio per il mio "M" o il tuo "M", ma dobbiamo considerare dove spostare il tuo "D" e se ciò creerà altri problemi. Che ne pensi di creare una lista delle opzioni di spostamento possibili e delle loro conseguenze per ogni oggetto, in modo da poter valutare meglio le possibilità e prendere una decisione informata? Inoltre, dovremmo anche considerare il numero di penalità che abbiamo accumulato e cercare di minimizzare il rischio di accumularne altre. 
	  } 
	   } 
	   } 
	 & & \\ 
 

    \theutterance \stepcounter{utterance}  

    & & \multicolumn{4}{p{0.6\linewidth}}{\cellcolor[rgb]{0.9,0.9,0.9}{%
	\makecell[{{p{\linewidth}}}]{% 
	  \tt {\tiny [GM$\rangle$B]}  
	 Il tuo messaggio è stato inoltrato all'altro giocatore.\\ \tt \\ \tt \\ \tt You are currently playing round 15 of maximum 28.\\ \tt You have collectively accumulated 15 of 16 penalties.\\ \tt L'altro giocatore ha inviato questo messaggio:\\ \tt \\ \tt \\ \tt "Sì, sono d'accordo! Valutare le conseguenze di ogni opzione è fondamentale per evitare di creare altri problemi. Spostare il mio "O" dalla posizione (5,5) per fare spazio per il tuo "M" potrebbe essere una buona opzione, ma dobbiamo considerare dove spostare il mio "O" e se ciò creerà altri problemi. Spostare il mio "D" dalla posizione (6,6) per fare spazio per il tuo "M" potrebbe essere un'altra opzione, ma dobbiamo considerare se ciò creerà spazio sufficiente per il tuo "M" e se ci saranno altri oggetti che potrebbero ostacolare lo spostamento. Spostare il tuo "D" dalla posizione (2,5) potrebbe essere un'altra opzione per creare spazio per il mio "M" o il tuo "M", ma dobbiamo considerare dove spostare il tuo "D" e se ciò creerà altri problemi. Che ne pensi di creare una lista delle opzioni di spostamento possibili e delle loro conseguenze per ogni oggetto, in modo da poter valutare meglio le possibilità e prendere una decisione informata? Inoltre, dovremmo anche considerare il numero di penalità che abbiamo accumulato e cercare di minimizzare il rischio di accumularne altre."\\ \tt \\ \tt \\ \tt Qual è il tuo prossimo comando? 
	  } 
	   } 
	   } 
	 & \\ 
 

    \theutterance \stepcounter{utterance}  

    & & & \multicolumn{4}{p{0.6\linewidth}}{\cellcolor[rgb]{1,0.85,0.72}{%
	\makecell[{{p{\linewidth}}}]{% 
	  \tt {\tiny [GM$\langle$B]}  
	 COMUNICA: Sono d'accordo che creare una lista delle opzioni di spostamento possibili e delle loro conseguenze per ogni oggetto sia una buona idea. In questo modo, potremmo valutare meglio le possibilità e prendere una decisione informata. Considerando le opzioni che abbiamo discusso, potremmo creare una lista con le seguenti opzioni: \\ \tt 1. Spostare il mio "M" nella posizione (5,5), spostando il tuo "O" \\ \tt 2. Spostare il mio "M" nella posizione (6,6), spostando il tuo "D" \\ \tt 3. Spostare il tuo "M" nella posizione (4,3) \\ \tt 4. Spostare il mio "D" per creare spazio per il tuo "M" o il mio "M". \\ \tt Per ogni opzione, dovremmo considerare le conseguenze, come ad esempio gli oggetti che potrebbero essere spostati, le celle che potrebbero essere occupate e le possibili penalità. Inoltre, come hai detto, dovremmo anche considerare il numero di penalità che abbiamo accumulato e cercare di minimizzare il rischio di accumularne altre. Che ne pensi di iniziare a creare questa lista e di valutarla insieme? 
	  } 
	   } 
	   } 
	  \\ 
 

    \theutterance \stepcounter{utterance}  

    & & & \multicolumn{2}{p{0.3\linewidth}}{\cellcolor[rgb]{0.95,0.95,0.95}{%
	\makecell[{{p{\linewidth}}}]{% 
	  \tt {\tiny [GM$|$GM]}  
	 COMUNICA: Sono d'accordo che creare una lista delle opzioni di spostamento possibili e delle loro conseguenze per ogni oggetto sia una buona idea. In questo modo, potremmo valutare meglio le possibilità e prendere una decisione informata. Considerando le opzioni che abbiamo discusso, potremmo creare una lista con le seguenti opzioni: \\ \tt 1. Spostare il mio "M" nella posizione (5,5), spostando il tuo "O" \\ \tt 2. Spostare il mio "M" nella posizione (6,6), spostando il tuo "D" \\ \tt 3. Spostare il tuo "M" nella posizione (4,3) \\ \tt 4. Spostare il mio "D" per creare spazio per il tuo "M" o il mio "M". \\ \tt Per ogni opzione, dovremmo considerare le conseguenze, come ad esempio gli oggetti che potrebbero essere spostati, le celle che potrebbero essere occupate e le possibili penalità. Inoltre, come hai detto, dovremmo anche considerare il numero di penalità che abbiamo accumulato e cercare di minimizzare il rischio di accumularne altre. Che ne pensi di iniziare a creare questa lista e di valutarla insieme? 
	  } 
	   } 
	   } 
	 & & \\ 
 

    \theutterance \stepcounter{utterance}  

    & & & \multicolumn{2}{p{0.3\linewidth}}{\cellcolor[rgb]{0.95,0.95,0.95}{%
	\makecell[{{p{\linewidth}}}]{% 
	  \tt {\tiny [GM$|$GM]}  
	 Invalid format: tail is not empty: '\\ \tt 1. Spostare il mio "M" nella posizione (5,5), spostando il tuo "O" \\ \tt 2. Spostare il mio "M" nella posizione (6,6), spostando il tuo "D" \\ \tt 3. Spostare il tuo "M" nella posizione (4,3) \\ \tt 4. Spostare il mio "D" per creare spazio per il tuo "M" o il mio "M". \\ \tt Per ogni opzione, dovremmo considerare le conseguenze, come ad esempio gli oggetti che potrebbero essere spostati, le celle che potrebbero essere occupate e le possibili penalità. Inoltre, come hai detto, dovremmo anche considerare il numero di penalità che abbiamo accumulato e cercare di minimizzare il rischio di accumularne altre. Che ne pensi di iniziare a creare questa lista e di valutarla insieme? 
	  } 
	   } 
	   } 
	 & & \\ 
 

    \theutterance \stepcounter{utterance}  

    & & \multicolumn{4}{p{0.6\linewidth}}{\cellcolor[rgb]{0.9,0.9,0.9}{%
	\makecell[{{p{\linewidth}}}]{% 
	  \tt {\tiny [GM$\rangle$B]}  
	 Penalità: Il tuo messaggio non deve contenere nulla dopo il comando!\\ \tt \\ \tt Assicurati che la tua risposta contenga solo `COMUNICA: <MESSAGE>` o `SPOSTA: <OBJECT>, (<X>, <Y>)`, e nient'altro!\\ \tt \\ \tt You have collectively accumulated 16 of 16 penalties.\\ \tt Riprova! 
	  } 
	   } 
	   } 
	 & \\ 
 

    \theutterance \stepcounter{utterance}  

    & & & \multicolumn{4}{p{0.6\linewidth}}{\cellcolor[rgb]{1,0.85,0.72}{%
	\makecell[{{p{\linewidth}}}]{% 
	  \tt {\tiny [GM$\langle$B]}  
	 COMUNICA: Sono d'accordo che creare una lista delle opzioni di spostamento possibili e delle loro conseguenze per ogni oggetto sia una buona idea, in questo modo potremmo valutare meglio le possibilità e prendere una decisione informata. 
	  } 
	   } 
	   } 
	  \\ 
 

    \theutterance \stepcounter{utterance}  

    & & & \multicolumn{2}{p{0.3\linewidth}}{\cellcolor[rgb]{0.95,0.95,0.95}{%
	\makecell[{{p{\linewidth}}}]{% 
	  \tt {\tiny [GM$|$GM]}  
	 COMUNICA: Sono d'accordo che creare una lista delle opzioni di spostamento possibili e delle loro conseguenze per ogni oggetto sia una buona idea, in questo modo potremmo valutare meglio le possibilità e prendere una decisione informata. 
	  } 
	   } 
	   } 
	 & & \\ 
 

    \theutterance \stepcounter{utterance}  

    & & \multicolumn{4}{p{0.6\linewidth}}{\cellcolor[rgb]{0.9,0.9,0.9}{%
	\makecell[{{p{\linewidth}}}]{% 
	  \tt {\tiny [A$\langle$GM]}  
	 Il tuo messaggio è stato inoltrato all'altro giocatore.\\ \tt \\ \tt \\ \tt You are currently playing round 15 of maximum 28.\\ \tt You have collectively accumulated 16 of 16 penalties.\\ \tt L'altro giocatore ha inviato questo messaggio:\\ \tt \\ \tt \\ \tt "Sono d'accordo che creare una lista delle opzioni di spostamento possibili e delle loro conseguenze per ogni oggetto sia una buona idea, in questo modo potremmo valutare meglio le possibilità e prendere una decisione informata."\\ \tt \\ \tt \\ \tt Qual è il tuo prossimo comando? 
	  } 
	   } 
	   } 
	 & \\ 
 

    \theutterance \stepcounter{utterance}  

    & \multicolumn{4}{p{0.6\linewidth}}{\cellcolor[rgb]{0.8,1,0.9}{%
	\makecell[{{p{\linewidth}}}]{% 
	  \tt {\tiny [A$\rangle$GM]}  
	 COMUNICA: Sì, sono d'accordo! Purtroppo, sembra che abbiamo raggiunto il limite massimo di penalità. Non possiamo più eseguire mosse senza rischiare di perdere la partita. In questo caso, potremmo cercare di valutare la situazione attuale e vedere se possiamo trovare una configurazione che funzioni per entrambi, anche se non è perfetta. Potremmo anche discutere se vale la pena continuare a giocare o se è meglio fermarsi e considerare la partita come persa. Che ne pensi? Dovremmo provare a trovare una soluzione o accettare la sconfitta? COMUNICA: finito? 
	  } 
	   } 
	   } 
	 & & \\ 
 

    \theutterance \stepcounter{utterance}  

    & & & \multicolumn{2}{p{0.3\linewidth}}{\cellcolor[rgb]{0.95,0.95,0.95}{%
	\makecell[{{p{\linewidth}}}]{% 
	  \tt {\tiny [GM$|$GM]}  
	 COMUNICA: Sì, sono d'accordo! Purtroppo, sembra che abbiamo raggiunto il limite massimo di penalità. Non possiamo più eseguire mosse senza rischiare di perdere la partita. In questo caso, potremmo cercare di valutare la situazione attuale e vedere se possiamo trovare una configurazione che funzioni per entrambi, anche se non è perfetta. Potremmo anche discutere se vale la pena continuare a giocare o se è meglio fermarsi e considerare la partita come persa. Che ne pensi? Dovremmo provare a trovare una soluzione o accettare la sconfitta? COMUNICA: finito? 
	  } 
	   } 
	   } 
	 & & \\ 
 

    \theutterance \stepcounter{utterance}  

    & & & \multicolumn{2}{p{0.3\linewidth}}{\cellcolor[rgb]{0.95,0.95,0.95}{%
	\makecell[{{p{\linewidth}}}]{% 
	  \tt {\tiny [GM$|$GM]}  
	 Invalid format: head is not empty: 'COMUNICA: Sì, sono d'accordo! Purtroppo, sembra che abbiamo raggiunto il limite massimo di penalità. Non possiamo più eseguire mosse senza rischiare di perdere la partita. In questo caso, potremmo cercare di valutare la situazione attuale e vedere se possiamo trovare una configurazione che funzioni per entrambi, anche se non è perfetta. Potremmo anche discutere se vale la pena continuare a giocare o se è meglio fermarsi e considerare la partita come persa. Che ne pensi? Dovremmo provare a trovare una soluzione o accettare la sconfitta? ' 
	  } 
	   } 
	   } 
	 & & \\ 
 

    \theutterance \stepcounter{utterance}  

    & & & \multicolumn{2}{p{0.3\linewidth}}{\cellcolor[rgb]{0.95,0.95,0.95}{%
	\makecell[{{p{\linewidth}}}]{% 
	  \tt {\tiny [GM$|$GM]}  
	 Maximum number of penalties exceeded 
	  } 
	   } 
	   } 
	 & & \\ 
 

    \theutterance \stepcounter{utterance}  

    & & & \multicolumn{2}{p{0.3\linewidth}}{\cellcolor[rgb]{0.95,0.95,0.95}{%
	\makecell[{{p{\linewidth}}}]{% 
	  \tt {\tiny [GM$|$GM]}  
	 * success: False\\ \tt * lose: True\\ \tt * aborted: True\\ \tt -------\\ \tt * Shifts: 0.00\\ \tt * Max Shifts: 12.00\\ \tt * Min Shifts: 6.00\\ \tt * End Distance Sum: 29.45\\ \tt * Init Distance Sum: 30.87\\ \tt * Expected Distance Sum: 29.33\\ \tt * Penalties: 17.00\\ \tt * Max Penalties: 16.00\\ \tt * Rounds: 15.00\\ \tt * Max Rounds: 28.00\\ \tt * Object Count: 7.00\\ \tt  
	  } 
	   } 
	   } 
	 & & \\ 
 

    \theutterance \stepcounter{utterance}  

    & & & \multicolumn{2}{p{0.3\linewidth}}{\cellcolor[rgb]{0.95,0.95,0.95}{%
	\makecell[{{p{\linewidth}}}]{% 
	  \tt {\tiny [GM$|$GM]}  
	 * Main Score: 0.00\\ \tt \\ \tt -------\\ \tt * Distance Score: 0.00\\ \tt * Coverage Score: 0.03\\ \tt * Penalty Score: 0.00\\ \tt * Alternative Penalty Score: 0.00\\ \tt * Alternative Main Score: 0.00\\ \tt \\ \tt -------\\ \tt * Shifts: 0.00\\ \tt * Max Shifts: 12.00\\ \tt * Min Shifts: 6.00\\ \tt * End Distance Sum: 29.45\\ \tt * Init Distance Sum: 30.87\\ \tt * Expected Distance Sum: 29.33\\ \tt * Penalties: 17.00\\ \tt * Max Penalties: 16.00\\ \tt * Rounds: 15.00\\ \tt * Max Rounds: 28.00\\ \tt * Object Count: 7.00\\ \tt  
	  } 
	   } 
	   } 
	 & & \\ 
 

\end{supertabular}
}

\end{document}
